% 用户定稿图：保持原图所表达的拓扑、轨线与局部模型信息。
% 两种固定间距的排线用于第 1 节的阴影/无阴影图。
\makeatletter
\pgfdeclarepatternformonly{milnorhatchdown}
  {\pgfqpoint{0pt}{-.75pt}}
  {\pgfqpoint{6.75pt}{6pt}}
  {\pgfqpoint{6pt}{6pt}}{%
    \pgfsetcolor{\tikz@pattern@color}%
    \pgfsetlinewidth{.75pt}%
    \pgfpathmoveto{\pgfqpoint{0pt}{6pt}}%
    \pgfpathlineto{\pgfqpoint{6.75pt}{-.75pt}}%
    \pgfusepath{stroke}}
\pgfdeclarepatternformonly{milnorhatchup}
  {\pgfqpoint{0pt}{0pt}}
  {\pgfqpoint{6.75pt}{6.75pt}}
  {\pgfqpoint{6pt}{6pt}}{%
    \pgfsetcolor{\tikz@pattern@color}%
    \pgfsetlinewidth{.75pt}%
    \pgfpathmoveto{\pgfqpoint{0pt}{0pt}}%
    \pgfpathlineto{\pgfqpoint{6.75pt}{6.75pt}}%
    \pgfusepath{stroke}}
\makeatother

\newcommand{\FigOneShadow}{%
\resizebox{.92\linewidth}{!}{%
\begin{tikzpicture}[x=.75pt,y=.75pt,yscale=-1,xscale=1,
  line cap=round,line join=round]
  % 整个圆柱形配边及右端的遮蔽轮廓。
  \draw[line width=1.5pt]
    (101.9,180.89)--(557.92,180.83)
    ..controls(583.39,180.82)and(604.15,215.87)..(604.29,259.10)
    ..controls(604.42,302.33)and(583.88,337.38)..(558.41,337.38)
    --(102.40,337.45)
    ..controls(76.92,337.45)and(56.16,302.41)..(56.03,259.18)
    ..controls(55.89,215.95)and(76.43,180.90)..(101.90,180.89)
    ..controls(127.38,180.89)and(148.13,215.93)..(148.27,259.17)
    ..controls(148.41,302.40)and(127.87,337.45)..(102.40,337.45);
  \draw[dash pattern=on 4.5pt off 4.5pt]
    (558.41,337.38)..controls(533.31,337.38)and(512.97,302.36)..(512.97,259.16)
    ..controls(512.97,215.96)and(533.31,180.93)..(558.41,180.93);
  \draw[dash pattern=on 4.5pt off 4.5pt]
    (406.67,181.35)..controls(429.85,179.30)and(447.28,216.19)..(447.28,259.26)
    ..controls(447.28,302.33)and(429.85,339.20)..(408.16,337.24);

  % c：大阴影区；小孔使用相反方向的排线。
  \path[pattern=milnorhatchdown,pattern color=black]
    (102.40,337.45)
    ..controls(229.28,337.28)and(181.28,337.28)..(408.16,337.24)
    ..controls(390.67,331.87)and(319.16,302.24)..(293.16,311.24)
    ..controls(265.16,319.24)and(239.16,300.24)..(239.16,262.24)
    ..controls(239.16,229.24)and(267.66,202.74)..(291.16,214.24)
    ..controls(339.16,233.24)and(366.16,208.24)..(411.55,182.32)
    ..controls(180.28,181.28)and(287.28,182.28)..(101.90,180.89)--cycle;
  % 阴影区只补画内部边界；闭合路径左端不应出现竖直描边。
  \draw[line width=1.5pt]
    (408.16,337.24)
    ..controls(390.67,331.87)and(319.16,302.24)..(293.16,311.24)
    ..controls(265.16,319.24)and(239.16,300.24)..(239.16,262.24)
    ..controls(239.16,229.24)and(267.66,202.74)..(291.16,214.24)
    ..controls(339.16,233.24)and(366.16,208.24)..(411.55,182.32);
  \draw[pattern=milnorhatchup,pattern color=black,line width=1.5pt]
    (335.28,248.28)
    ..controls(349.28,244.28)and(356.28,236.28)..(369.28,235.28)
    ..controls(382.28,234.28)and(394.28,241.28)..(387.28,260.28)
    ..controls(380.28,279.28)and(354.28,277.28)..(336.28,277.28)
    ..controls(318.28,277.28)and(299.28,270.28)..(300.28,252.28)
    ..controls(301.28,234.28)and(321.28,252.28)..(335.28,248.28)--cycle;

  % 左端口保持无阴影，以显出其为边界截面。
  \fill[white]
    (102.40,181.00)..controls(127.49,181.00)and(147.84,216.02)..(147.84,259.22)
    ..controls(147.84,302.43)and(127.49,337.45)..(102.40,337.45)--cycle;
  % 覆盖阴影图案在闭合裁切边缘留下的周期性细痕。
  \draw[white,line width=2.2pt,line cap=butt]
    (102.40,182.20)--(102.40,336.25);
  \draw[line width=.75pt]
    (102.40,181.00)..controls(127.49,181.00)and(147.84,216.02)..(147.84,259.22)
    ..controls(147.84,302.43)and(127.49,337.45)..(102.40,337.45);

  \node at (220,372) {$c$（阴影）};
  \node at (510,372) {$c'$（无阴影）};
  \path[use as bounding box] (45,170) rectangle (615,390);
\end{tikzpicture}}}

\newcommand{\FigTwoThree}{%
\resizebox{.98\linewidth}{!}{%
\begin{tikzpicture}[x=.75pt,y=.75pt,yscale=-1,xscale=1,
  >=Stealth,line width=.75pt,line cap=round,line join=round]
  % 图 2（左）：用户绘制的整体曲面。
  \draw (121.21,74.11)..controls(144.18,82.79)and(170.18,84.79)..(198.07,73.11);
  \draw (121.21,74.11)..controls(150.18,61.79)and(171.18,65.79)..(198.07,73.11);
  \draw (63.61,171.10)..controls(63.61,123.93)and(88.20,84.42)..(121.21,74.11);
  \draw (121.93,266.07)..controls(88.49,255.47)and(63.61,215.14)..(63.61,167.00);
  \draw (198.07,73.11)..controls(228.87,84.69)and(251.63,125.55)..(251.63,174.20);
  \draw (251.63,171.60)..controls(251.63,217.65)and(229.14,256.31)..(198.80,266.91);
  \draw (121.93,266.07)..controls(138.18,269.79)and(147.18,304.79)..(111.92,314.74);
  \draw (198.80,266.91)..controls(174.91,275.52)and(181.18,305.79)..(201.18,309.79);
  \draw (60.18,409.79)..controls(60.18,363.93)and(82.19,325.44)..(111.92,314.74);
  \draw (118.21,507.17)..controls(84.77,496.57)and(59.89,456.24)..(59.89,408.10);
  \draw (201.18,309.79)..controls(231.98,321.36)and(254.74,362.23)..(254.74,410.87);
  \draw (254.74,410.87)..controls(254.74,457.13)and(230.39,495.87)..(197.72,505.90);
  \draw (118.21,507.17)..controls(137.18,515.79)and(176.83,516.52)..(197.72,505.90);
  \draw[dash pattern=on 4.5pt off 4.5pt]
    (118.21,507.17)..controls(140.18,498.79)and(172.18,498.79)..(195.07,506.17);

  % 两个内边界分支。
  \draw (112.72,171.29)..controls(112.72,135.11)and(133.64,105.79)..(159.45,105.79)
    ..controls(185.26,105.79)and(206.18,135.11)..(206.18,171.29);
  \draw (206.18,171.29)..controls(206.18,207.46)and(185.26,236.79)..(159.45,236.79)
    ..controls(133.64,236.79)and(112.72,207.46)..(112.72,171.29);
  \draw (109.72,405.29)..controls(109.72,369.11)and(130.64,339.79)..(156.45,339.79)
    ..controls(182.26,339.79)and(203.18,369.11)..(203.18,405.29);
  \draw (203.18,405.29)..controls(203.18,441.46)and(182.26,470.79)..(156.45,470.79)
    ..controls(130.64,470.79)and(109.72,441.46)..(109.72,405.29);

  % 原书的四个临界点与边界标记。
  \fill (159.45,105.79) circle (2.1pt) node[above right=2pt] {$p_4$};
  \fill (159.45,236.79) circle (2.1pt) node[below right=2pt] {$p_3$};
  \fill (156.45,339.79) circle (2.1pt) node[above right=2pt] {$p_2$};
  \fill (156.45,512.10) circle (2.1pt) node[below right=3pt] {$p_1$};
  \node[above=5pt] at (159.45,70) {$V_1$};
  \node[below=12pt] at (156.45,512.10) {$V_0$};
  \node[left=7pt] at (62,290) {$W$};

  % 图 3（右）：用户绘制的四个分块。
  \draw (473.21,20.11)..controls(496.18,28.79)and(522.18,30.79)..(550.07,19.11);
  \draw (473.21,20.11)..controls(502.18,7.79)and(523.18,11.79)..(550.07,19.11);
  \draw (415.61,117.10)..controls(415.61,69.93)and(440.20,30.42)..(473.21,20.11);
  \draw (550.07,19.11)..controls(580.87,30.69)and(603.63,71.55)..(603.63,120.20);
  \draw (464.72,117.29)..controls(464.72,81.11)and(485.64,51.79)..(511.45,51.79)
    ..controls(537.26,51.79)and(558.18,81.11)..(558.18,117.29);
  \draw (473.93,237.07)..controls(440.49,226.47)and(415.61,186.14)..(415.61,138.00);
  \draw (603.63,142.60)..controls(603.63,188.65)and(581.14,227.31)..(550.80,237.91);
  \draw (558.18,142.29)..controls(558.18,178.46)and(537.26,207.79)..(511.45,207.79)
    ..controls(485.64,207.79)and(464.72,178.46)..(464.72,142.29);
  \draw (473.93,237.07)..controls(482.05,238.93)and(484.30,250.24)..(484.18,260.69);
  \draw (550.80,237.91)..controls(538.86,242.22)and(534.45,251.94)..(535.53,261.08);
  \draw (484.18,260.69)..controls(493.74,271.48)and(519.98,273.49)..(535.53,261.08);
  \draw[dash pattern=on 4.5pt off 4.5pt]
    (484.18,260.69)..controls(506.74,253.48)and(511.98,251.58)..(535.53,261.08);
  \draw (485.86,286.88)..controls(501.86,278.20)and(521.50,276.98)..(533.22,289.47);
  \draw (485.86,286.88)..controls(495.92,297.93)and(479.55,309.76)..(461.92,314.74);
  \draw (533.22,289.47)..controls(534.45,298.92)and(541.18,307.79)..(551.18,309.79);
  \draw (410.18,409.79)..controls(410.18,363.93)and(432.19,325.44)..(461.92,314.74);
  \draw (551.18,309.79)..controls(581.98,321.36)and(604.74,362.23)..(604.74,410.87);
  \draw (459.72,405.29)..controls(459.72,369.11)and(480.64,339.79)..(506.45,339.79)
    ..controls(532.26,339.79)and(553.18,369.11)..(553.18,405.29);
  \draw (410.89,433.10)..controls(410.89,481.24)and(435.77,521.57)..(469.21,532.17);
  \draw (605.74,435.87)..controls(605.74,482.13)and(581.39,520.87)..(548.72,530.90);
  \draw (554.18,430.29)..controls(554.18,466.46)and(533.26,495.79)..(507.45,495.79)
    ..controls(481.64,495.79)and(460.72,466.46)..(460.72,430.29);
  \draw (469.21,532.17)..controls(488.18,540.79)and(527.83,541.52)..(548.72,530.90);
  \draw[dash pattern=on 4.5pt off 4.5pt]
    (469.21,532.17)..controls(491.18,523.79)and(523.18,523.79)..(546.07,531.17);

  % 各分块端口的可见/遮蔽弧线。
  \draw (415.61,138.00)..controls(425.18,148.79)and(449.18,154.69)..(464.72,142.29);
  \draw (415.61,138.00)..controls(438.18,130.79)and(441.18,132.79)..(464.72,142.29);
  \draw (558.18,142.29)..controls(568.18,153.79)and(588.08,155.01)..(603.63,142.60);
  \draw (558.18,142.29)..controls(580.74,135.07)and(580.08,133.10)..(603.63,142.60);
  \draw (410.89,433.10)..controls(429.18,441.69)and(445.18,442.69)..(460.72,430.29);
  \draw (410.89,433.10)..controls(420.18,418.60)and(449.01,417.79)..(460.72,430.29);
  \draw (554.18,430.29)..controls(566.18,443.69)and(590.19,448.28)..(605.74,435.87);
  \draw (554.18,430.29)..controls(570.18,421.60)and(594.02,423.37)..(605.74,435.87);
  \draw (415.61,117.10)..controls(428.18,125.79)and(449.18,129.69)..(464.72,117.29);
  \draw (558.18,117.29)..controls(565.18,128.79)and(588.08,132.61)..(603.63,120.20);
  \draw[dash pattern=on 4.5pt off 4.5pt]
    (415.61,117.10)..controls(438.18,109.89)and(441.18,107.79)..(464.72,117.29);
  \draw[dash pattern=on 4.5pt off 4.5pt]
    (558.18,117.29)..controls(580.74,110.07)and(580.08,110.70)..(603.63,120.20);
  \draw (409.89,408.10)..controls(428.18,416.69)and(444.18,417.69)..(459.72,405.29);
  \draw (553.18,405.29)..controls(565.18,418.69)and(589.19,423.28)..(604.74,410.87);
  \draw[dash pattern=on 4.5pt off 4.5pt]
    (409.89,408.10)..controls(419.18,393.60)and(448.01,392.79)..(459.72,405.29);
  \draw[dash pattern=on 4.5pt off 4.5pt]
    (553.18,405.29)..controls(569.18,396.60)and(593.02,398.37)..(604.74,410.87);
  \draw (485.86,286.88)..controls(495.86,298.38)and(517.67,301.88)..(533.22,289.47);

  % 原图置于两图之间的坐标轴。
  \coordinate (axisorigin) at (330,620);
  \draw[->,line width=1pt] (axisorigin)--(330,552) node[above] {$z$};
  \draw[->,line width=1pt] (axisorigin)--(400,620) node[right] {$x$};
  \draw[->,line width=1pt] (axisorigin)--(286,670) node[below] {$y$};
  \node at (156,590) {图 2};
  \node at (508,590) {图 3};
  \path[use as bounding box] (12,0) rectangle (620,688);
\end{tikzpicture}}}

\newcommand{\FigThreeOne}{%
\begin{tikzpicture}[line cap=round,line join=round,scale=1.18]
  % 六个水平流形依次排列；中间两条分别位于临界水平两侧。
  \draw[thick] (-3.05,-1.75)..controls(-3.22,-1.15)and(-2.88,-.55)..(-3.05,.05)
    ..controls(-3.18,.60)and(-2.92,1.18)..(-3.05,1.75);
  \draw[thick] (-1.78,-1.75)..controls(-1.60,-1.20)and(-1.96,-.58)..(-1.78,.02)
    ..controls(-1.62,.62)and(-1.94,1.18)..(-1.78,1.75);
  \draw[thick] (-.36,-1.75)..controls(-.50,-1.12)and(-.20,-.55)..(-.36,.02)
    ..controls(-.50,.58)and(-.22,1.17)..(-.36,1.75);
  \draw[thick] (.36,-1.75)..controls(.20,-1.12)and(.50,-.55)..(.36,.02)
    ..controls(.20,.58)and(.50,1.17)..(.36,1.75);
  \draw[thick] (1.78,-1.75)..controls(1.94,-1.15)and(1.62,-.56)..(1.78,.04)
    ..controls(1.94,.62)and(1.60,1.18)..(1.78,1.75);
  \draw[thick] (3.05,-1.75)..controls(2.88,-1.15)and(3.22,-.55)..(3.05,.05)
    ..controls(2.90,.62)and(3.20,1.18)..(3.05,1.75);

  % U 同时覆盖 V_{-epsilon}、V_epsilon，p 位于两者之间且在 U 内。
  \draw[thick] (0,.02) ellipse (1.02 and .72);
  \fill (0,.02) circle (1.8pt) node[above=4pt] {$p$};
  \node[right=3pt] at (.93,.36) {$U$};

  % 六个水平流形的标记统一置于曲线下端之外，避免与线条叠压。
  \node[below=6pt] at (-3.05,-1.75) {$V$};
  \node[below=6pt] at (-1.78,-1.75) {$V_0$};
  \node[below=6pt,anchor=north east] at (-.10,-1.75) {$V_{-\epsilon}$};
  \node[below=6pt,anchor=north west] at (.10,-1.75) {$V_{\epsilon}$};
  \node[below=6pt] at (1.78,-1.75) {$V_1$};
  \node[below=6pt] at (3.05,-1.75) {$V'$};
\end{tikzpicture}}

\newcommand{\FigThreeTwo}{%
\resizebox{.96\linewidth}{!}{%
\begin{tikzpicture}[x=.75pt,y=.75pt,yscale=-1,xscale=1,
  >=Stealth,line width=.75pt,line cap=round,line join=round]
  % 用户绘制的二维指标 1 初等配边外形。
  \draw (150.72,141.79)..controls(183.72,137.79)and(241.97,141.41)..(303.72,125.79)
    ..controls(365.47,110.16)and(433.22,79.79)..(492.72,83.79)
    ..controls(552.22,87.79)and(561.72,198.79)..(518.72,210.79);
  \draw (132.72,361.79)..controls(112.22,329.79)and(95.10,287.04)..(99.72,243.79)
    ..controls(104.35,200.54)and(130.72,156.79)..(150.72,141.79);
  \draw[dash pattern=on 4.5pt off 4.5pt]
    (150.72,141.79)..controls(183.72,158.79)and(189.72,240.79)..(182.72,273.79)
    ..controls(175.72,306.79)and(159.72,335.79)..(132.72,361.79);
  \draw (132.72,361.79)..controls(216.72,350.79)and(236.72,357.79)..(296.72,377.79)
    ..controls(356.72,397.79)and(412.72,416.79)..(472.72,419.79);
  \draw (518.72,210.79)..controls(473.72,220.79)and(451.72,98.79)..(492.72,83.79);
  \draw (472.72,419.79)..controls(431.72,402.79)and(452.72,299.79)..(500.72,303.79)
    ..controls(548.72,307.79)and(527.72,433.79)..(472.72,419.79)--cycle;
  \draw (518.72,210.79)..controls(477.22,208.29)and(408.72,211.79)..(395.72,238.79)
    ..controls(382.72,265.79)and(389.72,307.79)..(500.72,303.79);
  \draw (452.72,352.79)..controls(341.72,329.79)and(299.72,293.79)..(334.72,219.79)
    ..controls(369.72,145.79)and(446.72,144.79)..(470.72,142.79);
  \draw (102.00,226.00)..controls(128.72,229.79)and(239.72,208.79)..(282.72,213.79)
    ..controls(325.72,218.79)and(374.72,237.79)..(390.72,250.79);
  \draw[dash pattern=on 4.5pt off 4.5pt]
    (183.72,260.79)..controls(204.22,262.79)and(228.72,270.79)..(278.72,271.79)
    ..controls(328.72,272.79)and(370.72,265.79)..(390.72,250.79);

  % 左、右手球面及两张特征圆盘上的方向。
  \draw[<->,line width=1pt]
    (171.92,264.29)..controls(121.40,285.81)and(58.13,329.47)..(29.00,314.00)
    ..controls(-.72,298.21)and(3.18,263.36)..(21.72,252.79)
    ..controls(39.71,242.53)and(58.56,228.65)..(91.72,226.79);
  \draw[->,line width=1pt] (241.72,183.79)--(239.00,210.00);
  \draw[->,line width=1pt] (290.00,309.00)--(321.72,304.79);
  \draw[<->,line width=1pt]
    (491.28,142.19)..controls(530.38,143.11)and(623.07,156.25)..(616.28,255.19)
    ..controls(609.42,355.15)and(515.16,371.55)..(471.28,356.19);

  % 用户图中的交点；p 已由用户标出，不重复增设。
  \foreach \x/\y in {333.10/224.10,104.34/225.62,186.05/260.40,470.72/142.79,455.06/352.40}
    \fill (\x,\y) circle (2.35pt);
  \node[above left=3pt] at (333.10,224.10) {$p$};
  \node[above=4pt] at (225,126) {$W$};
  \node[anchor=east] at (-5,286) {$S_L$};
  \node[above=3pt] at (239,183.79) {$D_L$};
  \node[below=3pt] at (290,309) {$D_R$};
  \node[right=3pt] at (616.28,255.19) {$S_R$};
  \node[below=7pt] at (114,382) {$V$};
  \node[below=7pt] at (493,430) {$V'$};
  \path[use as bounding box] (-35,65) rectangle (640,456);
\end{tikzpicture}}}

% 图 3.3 与图 3.4 共用的沙漏形外轮廓。
\newcommand{\MilnorThreeHourglassOutline}{%
  \draw (542.22,89.21)..controls(497.70,153.04)and(419.38,195.40)..(330.20,195.40)
    ..controls(241.02,195.40)and(162.70,153.04)..(118.18,89.21);
  \draw (118.18,500.19)..controls(162.70,436.36)and(241.02,394.00)..(330.20,394.00)
    ..controls(419.38,394.00)and(497.70,436.36)..(542.22,500.19);
  \draw (83.21,129.59)..controls(120.23,174.19)and(142.80,233.50)..(142.80,298.60)
    ..controls(142.80,363.70)and(120.23,423.01)..(83.21,467.61);
  \draw (570.01,459.44)..controls(533.77,414.20)and(512.24,354.50)..(513.38,289.41)
    ..controls(514.51,224.33)and(538.11,165.42)..(575.91,121.47);}

\newcommand{\FigThreeThree}{%
\resizebox{.84\linewidth}{!}{%
\begin{tikzpicture}[x=.75pt,y=.75pt,yscale=-1,xscale=1,
  >=Stealth,line width=.75pt,line cap=round,line join=round]
  \MilnorThreeHourglassOutline

  % 用户图中的坐标轴、C 的圆柱形邻域和 D_L。
  \draw[->] (49.11,297.10)--(609.31,297.10);
  \draw[->] (329.31,539.10)--(329.31,67.50);
  \draw (139.94,258.80)--(516.20,253.60);
  \draw (139.94,337.40)--(516.20,333.21);
  \draw (139.94,258.80)..controls(141.82,271.76)and(142.80,285.05)..(142.80,298.60)
    ..controls(142.80,312.15)and(141.82,325.44)..(139.94,338.40);
  \draw (516.20,333.21)..controls(514.32,320.24)and(513.34,306.95)..(513.34,293.40)
    ..controls(513.34,279.86)and(514.32,266.56)..(516.20,253.60);
  \draw[pattern=milnorhatchup,pattern color=black]
    (141.58,259.13)--(513.53,253.28)--(514.77,332.37)--(142.82,338.22)--cycle;
  \draw[line width=2.25pt] (142.20,298.80)--(514.20,295.80);

  % 用户已画出左上四分之一；按横、纵中心轴反射，补齐其余三部分。
  \foreach \xs/\ys in {1/1,-1/1,1/-1,-1/-1}{
    \begin{scope}[shift={(329.31,297.10)},xscale=\xs,yscale=\ys]
      \draw[->] (-196.11,-174.50)--(-222.11,-149.50);
      \draw[->] (-173.11,-149.50)--(-199.11,-124.50);
      \draw[->] (-151.11,-124.50)--(-177.11,-99.50);
      \draw[->] (-127.11,-105.50)--(-150.11,-75.50);
      \draw[->] (-96.11,-92.50)--(-102.11,-59.50);
      \draw[->] (-55.11,-86.50)--(-55.11,-55.50);
    \end{scope}}

  % D_L、C 的指示箭头与两侧的 V 标记。
  \draw[->] (96.20,273.80)--(138.00,292.20);
  \draw[->] (555.20,247.80)--(500.20,284.80);
  \node at (137,458) {$V$};
  \node at (520,458) {$V$};
  \node at (74,258) {$D_L$};
  \node at (578,240) {$C$};
  \path[use as bounding box] (40,55) rectangle (620,525);
\end{tikzpicture}}}

\newcommand{\FigThreeFour}{%
\resizebox{.84\linewidth}{!}{%
\begin{tikzpicture}[x=.75pt,y=.75pt,yscale=-1,xscale=1,
  >=Stealth,line width=.75pt,line cap=round,line join=round]
  \MilnorThreeHourglassOutline

  % 图 3.4 复用图 3.3 的轮廓；中央矩形表示 C 中的第二次收缩。
  \draw (142.80,258.80)--(513.34,258.80);
  \draw (142.80,337.40)--(513.34,337.40);
  \draw (142.80,258.80)--(142.80,337.40);
  \draw (513.34,258.80)--(513.34,337.40);
  \draw[line width=1.15pt] (142.80,297.10)--(513.34,297.10);

  % 情形 2：上下两块沿直线移向 V union D_L。
  \foreach \x in {180,225,270,329,388,433,478}{
    \draw[->] (\x,211)--(\x,255);
    \draw[->] (\x,382)--(\x,341);}

  % 情形 1 与情形 3：靠近两侧边界时分别沿侧向轨线收缩。
  \draw[->] (120,205)..controls(132,220)and(138,237)..(141,252);
  \draw[->] (538,205)..controls(526,220)and(519,237)..(515,252);
  \draw[->] (120,389)..controls(132,373)and(138,354)..(141,341);
  \draw[->] (538,389)..controls(526,373)and(519,354)..(515,341);

  \node[anchor=east] at (104,175) {情形 1};
  \node[anchor=west] at (555,175) {情形 1};
  \node at (329,282) {情形 2};
  \node[anchor=east] at (104,418) {情形 3};
  \node[anchor=west] at (555,418) {情形 3};
  \path[use as bounding box] (40,65) rectangle (620,510);
\end{tikzpicture}}}

\newcommand{\FigFourOne}{%
\begin{tikzpicture}[>=Stealth,x=1.35cm,y=1.08cm,
    line cap=round,line join=round]
  % 两侧水平集；各分叉的四个脚都是曲线上的显式节点。
  \coordinate (v0ta) at (-2.35,1.35);
  \coordinate (v0tb) at (-2.45,.55);
  \coordinate (v0ba) at (-2.43,-.45);
  \coordinate (v0bb) at (-2.32,-1.30);
  \coordinate (v1ta) at (2.38,1.30);
  \coordinate (v1tb) at (2.32,.60);
  \coordinate (v1ba) at (2.30,-.48);
  \coordinate (v1bb) at (2.46,-1.28);
  \draw[thick] (-2.28,1.74)
    .. controls (-2.31,1.58) and (-2.34,1.46) .. (v0ta)
    .. controls (-2.40,1.08) and (-2.43,.78) .. (v0tb)
    .. controls (-2.50,.16) and (-2.48,-.10) .. (v0ba)
    .. controls (-2.39,-.78) and (-2.35,-1.06) .. (v0bb)
    .. controls (-2.29,-1.48) and (-2.26,-1.62) .. (-2.25,-1.74);
  \draw[thick] (2.31,1.74)
    .. controls (2.34,1.57) and (2.36,1.44) .. (v1ta)
    .. controls (2.37,1.02) and (2.34,.77) .. (v1tb)
    .. controls (2.26,.15) and (2.27,-.12) .. (v1ba)
    .. controls (2.35,-.76) and (2.42,-1.05) .. (v1bb)
    .. controls (2.50,-1.48) and (2.53,-1.62) .. (2.55,-1.74);

  % 上、下两个圆滑 X 型分叉。
  \coordinate (p) at (-.55,.94);
  \draw[thick] (v0ta)
    .. controls (-1.72,1.31) and (-1.02,1.12) .. (p)
    .. controls (.15,.72) and (1.35,.64) .. (v1tb);
  \draw[thick] (v0tb)
    .. controls (-1.72,.58) and (-1.03,.70) .. (p)
    .. controls (.16,1.16) and (1.34,1.24) .. (v1ta);
  \coordinate (pp) at (.18,-.91);
  \draw[thick] (v0ba)
    .. controls (-1.62,-.46) and (-.55,-.61) .. (pp)
    .. controls (.84,-1.18) and (1.63,-1.27) .. (v1bb);
  \draw[thick] (v0bb)
    .. controls (-1.58,-1.22) and (-.55,-1.13) .. (pp)
    .. controls (.86,-.62) and (1.61,-.52) .. (v1ba);
  \fill (p) circle (1.7pt);
  \fill (pp) circle (1.7pt);
  \node[above=5pt] at (p) {$p$};
  \node[below=5pt] at (pp) {$p'$};
  \node[anchor=west] at (2.68,.94) {$K_p$};
  \node[anchor=west] at (2.68,-.91) {$K_{p'}$};
  \node[below=5pt] at (-2.28,-1.74) {$V_0$};
  \node[below=5pt] at (2.55,-1.74) {$V_1$};
\end{tikzpicture}}

\newcommand{\FigFourTwo}{%
\begin{tikzpicture}[>=Stealth,x=6.65cm,y=6.65cm,
    line cap=round,line join=round]
  \draw[->] (0,0)--(1.12,0) node[right] {$x$}; \draw[->] (0,0)--(0,1.12) node[above] {$z$};
  \draw[gray] (0,1)--(1,1); \draw[gray] (1,0)--(1,1);
  % 两条曲线在 x=0,1 的邻域内都严格等于 z=x；在指定临界值
  % 附近留出斜率为 1 的直线段，以直观对应条件 (iii)。
  \draw[thick]
    (0,0)--(.08,.08)
    .. controls (.13,.13) and (.18,.52) .. (.23,.57)
    --(.33,.67)
    .. controls (.40,.74) and (.85,.85) .. (.92,.92)
    --(1,1);
  \draw[thick]
    (0,0)--(.08,.08)
    .. controls (.15,.15) and (.60,.26) .. (.67,.33)
    --(.77,.43)
    .. controls (.82,.48) and (.87,.87) .. (.92,.92)
    --(1,1);
  % a=G(f(p),0), a'=G(f(p'),1)。
  \draw[dashed] (.28,0)--(.28,.62);
  \draw[dashed] (.72,0)--(.72,.38);
  \node[below] at (.28,0) {$f(p)$}; \node[below] at (.72,0) {$f(p')$};
  \node[left] at (0,.62) {$a$}; \node[left] at (0,.38) {$a'$};
  \node[above=5pt] at (.56,.80) {$z=G(x,0)$};
  \node[below=8pt] at (.56,.20) {$z=G(x,1)$};
\end{tikzpicture}}

\newcommand{\FigFourFour}{%
\begin{tikzpicture}[>=Stealth,x=.52pt,y=.52pt,yscale=-1,
    line width=.65pt,line cap=round,line join=round]
  % 左右边界以及两个中间水平集。
  \draw (122.2,503.6) .. controls (156.2,388.6) and (104.45,369.1)
    .. (93.2,283.6) .. controls (81.95,198.1) and (117.2,142.6) .. (168.2,118.6);
  \draw (287.2,502.6) .. controls (321.2,387.6) and (288.45,388.1)
    .. (277.2,302.6) .. controls (265.95,217.1) and (281.2,180.6) .. (313.2,115.6);
  \draw (352.2,501.6) .. controls (386.2,386.6) and (351.45,390.1)
    .. (340.2,304.6) .. controls (328.95,219.1) and (344.2,182.6) .. (376.2,117.6);
  \draw (592.2,498.6) .. controls (539.2,377.6) and (542.2,368.6)
    .. (550.2,279.6) .. controls (558.2,190.6) and (557.2,194.6) .. (531.2,122.6);

  % f^{-1}(a) 与 V 之间的轨线。
  \draw (288.14,170.43) .. controls (312.34,173.03) and (319.14,152.43) .. (356.14,161.43);
  \draw (279.2,198.6) .. controls (303.4,201.2) and (310.14,176.43) .. (347.14,185.43);
  \draw (275.2,224.6) .. controls (299.4,227.2) and (304,203) .. (341,212);
  \draw (273.2,255.6) .. controls (297.4,258.2) and (300.2,227.6) .. (337.2,236.6);
  \draw[dashed] (273.2,255.6) .. controls (300.2,264.6) and (315,262) .. (336,260);
  \draw[dashed] (278.2,312.6) .. controls (319.2,327.6) and (322.2,317.6) .. (342.2,313.6);
  \draw (278.2,312.6) .. controls (302.4,315.2) and (301.2,282.6) .. (338.2,291.6);

  % 经过 p 与 p' 的两组轨线。
  \draw (278.2,312.6) .. controls (254.2,311.6) and (249.2,312.6)
    .. (238.2,311.6) .. controls (217.2,310.6) and (197.2,303.6)
    .. (196.2,287.6) .. controls (195.2,272.6) and (208.2,262.6)
    .. (221.2,258.6) .. controls (241.2,252.6) and (252.2,255.6) .. (273.2,255.6);
  \draw (91,263) .. controls (94.2,269.6) and (99.2,272.6)
    .. (126.2,272.6) .. controls (153.2,272.6) and (162.2,270.6)
    .. (180.2,273.6) .. controls (192.2,275.6) and (196.2,280.6)
    .. (196.2,287.6) .. controls (196.7,297.6) and (164.2,306.6)
    .. (144.2,307.6) .. controls (124.2,308.6) and (112.2,312.6) .. (104.2,324.6);
  \draw (342.2,313.6) .. controls (384.2,300.6) and (436.2,319.6) .. (444.2,343.6);
  \draw (361.2,386.6) .. controls (408.2,384.6) and (438.2,385.6) .. (444.2,343.6);
  \draw (444.2,343.6) .. controls (484.2,313.6) and (507.2,317.6) .. (547.2,317.6);
  \draw (444.2,343.6) .. controls (459.4,372.2) and (514.2,379.6) .. (548.2,375.6);

  % 同痕造成的变形方向。
  \draw[<-] (350.15,231.5) .. controls (387.77,206.05) and (399.2,252.85) .. (438.2,223.6);
  \draw[<-] (354.01,296.22) .. controls (436.04,308.77) and (424.41,262.02) .. (438.2,223.6);
  \draw[<-] (345.98,262.8) .. controls (372.86,274.15) and (436.6,294.55) .. (475.2,265.6);
  \draw[<-] (349.58,308.07) .. controls (443.94,296.56) and (435.8,295.15) .. (475.2,265.6);
  \draw[->] (264.2,396.6) .. controls (303.4,367.2) and (298.42,352.21) .. (338.2,321.6);
  \draw[->] (264.2,396.6) .. controls (285.96,413.85) and (312.13,421.3) .. (353.2,391.6);

  \foreach \x/\y in {342.2/313.6,341/291.6,333.2/260,337.2/239.4,
      444.2/343.6,196.2/287.6,361.2/386.6}
    {\fill (\x,\y) circle[radius=1.45pt];}

  \node[below] at (285,502.6) {$f^{-1}(a)$};
  \node[below] at (352.2,501.6) {$V$};
  \node[anchor=west] at (486,265.6) {$S_R$};
  \node[anchor=west] at (441,211) {$\overline{S_R}=h(S_R)$};
  \node[below left=2pt] at (196.2,287.6) {$p$};
  \node[below=4pt] at (444.2,343.6) {$p'$};
  \node[anchor=east] at (264.2,396.6) {$S'_L$};
  \path[use as bounding box] (75,105) rectangle (620,540);
\end{tikzpicture}}

\newcommand{\FigFiveOne}{%
\begin{tikzpicture}[x=.75pt,y=.75pt,yscale=-1,
    line width=.7pt,line cap=round,line join=round]
  % 以用户底稿为准：二维消去所给出的折返柄体。
  \draw (100,181.8)
    .. controls (100,142.7) and (112.58,111) .. (128.1,111)
    .. controls (143.62,111) and (156.2,142.7) .. (156.2,181.8)
    .. controls (156.2,220.9) and (143.62,252.6) .. (128.1,252.6)
    .. controls (112.58,252.6) and (100,220.9) .. (100,181.8)--cycle;
  \draw (128.1,111)
    .. controls (151.15,118.3) and (259.2,113.6) .. (274.2,111.6)
    .. controls (289.2,109.6) and (292.2,106.6) .. (307.2,106.6)
    .. controls (322.2,106.6) and (337.2,140.6) .. (311.2,147.6)
    .. controls (285.2,154.6) and (301.2,150.6) .. (285.2,155.6)
    .. controls (269.2,160.6) and (254.2,175.6) .. (250.2,180.6)
    .. controls (246.2,185.6) and (264.2,198.6) .. (291.2,201.6)
    .. controls (318.2,204.6) and (364.2,198.6) .. (371.2,196.6)
    .. controls (378.2,194.6) and (388.2,198.6) .. (392.2,218.6)
    .. controls (396.2,238.6) and (386.2,249.6) .. (381.2,250.6)
    .. controls (376.2,251.6) and (340.2,256.6) .. (286.2,250.6)
    .. controls (232.2,244.6) and (163.1,246.6) .. (128.1,252.6);
  \draw[dashed] (371.2,196.6)
    .. controls (350.2,197.6) and (354.2,256.6) .. (381.2,250.6);
  \draw (286.2,250.6)
    .. controls (293.2,236.6) and (297.2,216.6) .. (291.2,201.6);
  \draw (274.2,111.6)
    .. controls (284.2,123.6) and (284.2,136.6) .. (285.2,155.6);
  \path[use as bounding box] (88,96) rectangle (406,266);
\end{tikzpicture}}

\newcommand{\FigFiveTwo}{%
  \resizebox{!}{.52\textheight}{\tikzset{every picture/.style={line width=0.75pt}} %set default line width to 0.75pt        

\begin{tikzpicture}[x=0.75pt,y=0.75pt,yscale=-1,xscale=1]
%uncomment if require: \path (0,724); %set diagram left start at 0, and has height of 724

%Shape: Arc [id:dp771786952376691] 
\draw  [draw opacity=0] (121.27,634.98) .. controls (102.28,633.73) and (87,583.71) .. (87,522.2) .. controls (87,459.9) and (102.67,409.4) .. (122,409.4) -- (122,522.2) -- cycle ; \draw   (121.27,634.98) .. controls (102.28,633.73) and (87,583.71) .. (87,522.2) .. controls (87,459.9) and (102.67,409.4) .. (122,409.4) ;  
%Curve Lines [id:da8545165027851835] 
\draw    (122,409.4) .. controls (162.2,402) and (219.4,383.6) .. (245.2,383) .. controls (271,382.4) and (295.2,401) .. (297.2,432) .. controls (299.2,463) and (287.2,465) .. (277.2,484) .. controls (267.2,503) and (278.2,538) .. (317.2,545) .. controls (356.2,552) and (422.2,562) .. (464.2,567) ;
%Shape: Ellipse [id:dp817226227673999] 
\draw   (440.1,609.3) .. controls (440.1,585.94) and (450.89,567) .. (464.2,567) .. controls (477.51,567) and (488.3,585.94) .. (488.3,609.3) .. controls (488.3,632.66) and (477.51,651.6) .. (464.2,651.6) .. controls (450.89,651.6) and (440.1,632.66) .. (440.1,609.3) -- cycle ;
%Curve Lines [id:da9839094781377016] 
\draw    (121.27,634.98) .. controls (161.2,643) and (372.2,670) .. (464.2,651.6) ;
%Curve Lines [id:da6651840314711972] 
\draw    (97,443.4) .. controls (140.2,427) and (175.2,423) .. (212.2,416) .. controls (246.61,409.49) and (254.64,421.92) .. (249.56,433.35) ;
\draw [shift={(248.2,435.9)}, rotate = 302.01] [fill={rgb, 255:red, 0; green, 0; blue, 0 }  ][line width=0.08]  [draw opacity=0] (8.93,-4.29) -- (0,0) -- (8.93,4.29) -- cycle    ;
\draw [shift={(159.17,425.71)}, rotate = 168.1] [fill={rgb, 255:red, 0; green, 0; blue, 0 }  ][line width=0.08]  [draw opacity=0] (8.93,-4.29) -- (0,0) -- (8.93,4.29) -- cycle    ;
\draw [shift={(242.8,416.43)}, rotate = 195.59] [fill={rgb, 255:red, 0; green, 0; blue, 0 }  ][line width=0.08]  [draw opacity=0] (8.93,-4.29) -- (0,0) -- (8.93,4.29) -- cycle    ;
%Curve Lines [id:da1452296613462677] 
\draw    (273.33,444.64) .. controls (279.38,438.32) and (282.8,436.25) .. (292.7,443.9) ;
\draw [shift={(271.2,446.9)}, rotate = 313.15] [fill={rgb, 255:red, 0; green, 0; blue, 0 }  ][line width=0.08]  [draw opacity=0] (8.93,-4.29) -- (0,0) -- (8.93,4.29) -- cycle    ;
%Curve Lines [id:da7810401101193405] 
\draw    (263.2,389.9) .. controls (279.7,409.9) and (267.7,422.9) .. (265.7,430.4) ;
\draw [shift={(271.42,414.64)}, rotate = 272.83] [fill={rgb, 255:red, 0; green, 0; blue, 0 }  ][line width=0.08]  [draw opacity=0] (8.93,-4.29) -- (0,0) -- (8.93,4.29) -- cycle    ;
%Curve Lines [id:da3660635328346252] 
\draw    (88.2,507) .. controls (131.4,490.6) and (138.2,491) .. (175.2,484) .. controls (209.79,477.46) and (226.03,481.75) .. (227.59,494.6) ;
\draw [shift={(227.7,497.4)}, rotate = 271.91] [fill={rgb, 255:red, 0; green, 0; blue, 0 }  ][line width=0.08]  [draw opacity=0] (8.93,-4.29) -- (0,0) -- (8.93,4.29) -- cycle    ;
\draw [shift={(136.13,491.48)}, rotate = 166.21] [fill={rgb, 255:red, 0; green, 0; blue, 0 }  ][line width=0.08]  [draw opacity=0] (8.93,-4.29) -- (0,0) -- (8.93,4.29) -- cycle    ;
\draw [shift={(209.64,481.35)}, rotate = 183.53] [fill={rgb, 255:red, 0; green, 0; blue, 0 }  ][line width=0.08]  [draw opacity=0] (8.93,-4.29) -- (0,0) -- (8.93,4.29) -- cycle    ;
%Curve Lines [id:da6129744214808652] 
\draw    (249.56,502.04) .. controls (262.38,492.14) and (264.03,495.16) .. (269.7,501.4) ;
\draw [shift={(247.2,503.9)}, rotate = 321.11] [fill={rgb, 255:red, 0; green, 0; blue, 0 }  ][line width=0.08]  [draw opacity=0] (8.93,-4.29) -- (0,0) -- (8.93,4.29) -- cycle    ;
%Curve Lines [id:da8665971491876462] 
\draw    (265.7,430.4) .. controls (256.2,452.9) and (244.2,479.4) .. (240.2,495.4) ;
\draw [shift={(250.05,467.39)}, rotate = 292.34] [fill={rgb, 255:red, 0; green, 0; blue, 0 }  ][line width=0.08]  [draw opacity=0] (8.93,-4.29) -- (0,0) -- (8.93,4.29) -- cycle    ;
%Curve Lines [id:da3813998619722765] 
\draw    (240.2,495.4) .. controls (230.9,537.4) and (225.5,577.65) .. (255.41,595.81) .. controls (285.32,613.98) and (349.25,618.55) .. (441.2,620.2) ;
\draw [shift={(233.1,554.55)}, rotate = 267.75] [fill={rgb, 255:red, 0; green, 0; blue, 0 }  ][line width=0.08]  [draw opacity=0] (8.93,-4.29) -- (0,0) -- (8.93,4.29) -- cycle    ;
\draw [shift={(352.01,616.71)}, rotate = 184.4] [fill={rgb, 255:red, 0; green, 0; blue, 0 }  ][line width=0.08]  [draw opacity=0] (8.93,-4.29) -- (0,0) -- (8.93,4.29) -- cycle    ;
%Curve Lines [id:da25413126891176796] 
\draw    (119.2,416) .. controls (124.2,422) and (123.2,422) .. (127.2,429) ;
%Curve Lines [id:da9623357897940251] 
\draw    (133.2,441) .. controls (138.2,454) and (142.2,467) .. (144.2,483) ;
%Curve Lines [id:da2958692841491226] 
\draw    (145.2,498) .. controls (149.2,553) and (137.2,615) .. (121.27,634.98) ;
%Curve Lines [id:da028571924348444022] 
\draw  [dash pattern={on 4.5pt off 4.5pt}]  (254.2,402.9) .. controls (259.7,386.4) and (280.4,398.8) .. (285.2,417.4) .. controls (290,436) and (273.2,468.9) .. (269.7,479.9) .. controls (266.2,490.9) and (261.7,506.9) .. (256.7,517.4) .. controls (251.7,527.9) and (222.2,521.9) .. (220.7,502.9) .. controls (219.2,483.9) and (224.2,474.4) .. (235.2,447.9) .. controls (246.2,421.4) and (235.7,446.4) .. (254.2,402.9) -- cycle ;
%Shape: Arc [id:dp11600535545844648] 
\draw  [draw opacity=0] (120.27,309.12) .. controls (101.28,307.87) and (86,257.86) .. (86,196.35) .. controls (86,134.05) and (101.67,83.55) .. (121,83.55) -- (121,196.35) -- cycle ; \draw   (120.27,309.12) .. controls (101.28,307.87) and (86,257.86) .. (86,196.35) .. controls (86,134.05) and (101.67,83.55) .. (121,83.55) ;  
%Curve Lines [id:da8485262314692139] 
\draw    (121,83.55) .. controls (161.2,76.15) and (218.4,57.75) .. (244.2,57.15) .. controls (270,56.55) and (294.2,75.15) .. (296.2,106.15) .. controls (298.2,137.15) and (286.2,139.15) .. (276.2,158.15) .. controls (266.2,177.15) and (277.2,212.15) .. (316.2,219.15) .. controls (355.2,226.15) and (421.2,236.15) .. (463.2,241.15) ;
%Shape: Ellipse [id:dp27803344135238306] 
\draw   (439.1,283.45) .. controls (439.1,260.09) and (449.89,241.15) .. (463.2,241.15) .. controls (476.51,241.15) and (487.3,260.09) .. (487.3,283.45) .. controls (487.3,306.81) and (476.51,325.75) .. (463.2,325.75) .. controls (449.89,325.75) and (439.1,306.81) .. (439.1,283.45) -- cycle ;
%Curve Lines [id:da3705907763733033] 
\draw    (120.27,309.12) .. controls (160.2,317.15) and (371.2,344.15) .. (463.2,325.75) ;
%Curve Lines [id:da7989673700463524] 
\draw    (96,117.55) .. controls (139.2,101.15) and (174.2,97.15) .. (211.2,90.15) .. controls (248.2,83.15) and (257.2,88.15) .. (267.2,92.15) ;
\draw [shift={(158.17,99.86)}, rotate = 168.1] [fill={rgb, 255:red, 0; green, 0; blue, 0 }  ][line width=0.08]  [draw opacity=0] (8.93,-4.29) -- (0,0) -- (8.93,4.29) -- cycle    ;
\draw [shift={(244.59,86.69)}, rotate = 179.95] [fill={rgb, 255:red, 0; green, 0; blue, 0 }  ][line width=0.08]  [draw opacity=0] (8.93,-4.29) -- (0,0) -- (8.93,4.29) -- cycle    ;
%Shape: Circle [id:dp17199567700794005] 
\draw  [fill={rgb, 255:red, 0; green, 0; blue, 0 }  ,fill opacity=1 ] (262.4,92.15) .. controls (262.4,90.82) and (263.47,89.75) .. (264.8,89.75) .. controls (266.13,89.75) and (267.2,90.82) .. (267.2,92.15) .. controls (267.2,93.47) and (266.13,94.55) .. (264.8,94.55) .. controls (263.47,94.55) and (262.4,93.47) .. (262.4,92.15) -- cycle ;
%Curve Lines [id:da33191565010475343] 
\draw    (267.2,92.15) .. controls (279.2,98.15) and (286.2,105.15) .. (296.2,117.15) ;
\draw [shift={(278.16,98.88)}, rotate = 39.6] [fill={rgb, 255:red, 0; green, 0; blue, 0 }  ][line width=0.08]  [draw opacity=0] (8.93,-4.29) -- (0,0) -- (8.93,4.29) -- cycle    ;
%Curve Lines [id:da47257040187603183] 
\draw    (268.2,62.15) .. controls (269.4,87.75) and (270.2,68.15) .. (264.8,94.55) ;
\draw [shift={(267.34,83.5)}, rotate = 286.11] [fill={rgb, 255:red, 0; green, 0; blue, 0 }  ][line width=0.08]  [draw opacity=0] (8.93,-4.29) -- (0,0) -- (8.93,4.29) -- cycle    ;
%Curve Lines [id:da917490869937732] 
\draw    (87.2,181.15) .. controls (130.4,164.75) and (137.2,165.15) .. (174.2,158.15) .. controls (211.2,151.15) and (229.2,163.15) .. (239.2,167.15) ;
\draw [shift={(135.13,165.63)}, rotate = 166.21] [fill={rgb, 255:red, 0; green, 0; blue, 0 }  ][line width=0.08]  [draw opacity=0] (8.93,-4.29) -- (0,0) -- (8.93,4.29) -- cycle    ;
\draw [shift={(212.6,157.58)}, rotate = 188.93] [fill={rgb, 255:red, 0; green, 0; blue, 0 }  ][line width=0.08]  [draw opacity=0] (8.93,-4.29) -- (0,0) -- (8.93,4.29) -- cycle    ;
%Shape: Circle [id:dp7405970458433607] 
\draw  [fill={rgb, 255:red, 0; green, 0; blue, 0 }  ,fill opacity=1 ] (236.8,167.15) .. controls (236.8,165.82) and (237.87,164.75) .. (239.2,164.75) .. controls (240.53,164.75) and (241.6,165.82) .. (241.6,167.15) .. controls (241.6,168.47) and (240.53,169.55) .. (239.2,169.55) .. controls (237.87,169.55) and (236.8,168.47) .. (236.8,167.15) -- cycle ;
%Curve Lines [id:da07557205314940973] 
\draw    (239.2,167.15) .. controls (249.2,174.15) and (259.2,177.15) .. (273.2,179.15) ;
\draw [shift={(249.4,172.95)}, rotate = 21.22] [fill={rgb, 255:red, 0; green, 0; blue, 0 }  ][line width=0.08]  [draw opacity=0] (8.93,-4.29) -- (0,0) -- (8.93,4.29) -- cycle    ;
%Curve Lines [id:da12569593671542] 
\draw    (239.2,167.15) .. controls (245.4,145.75) and (252.2,123.15) .. (264.8,92.15) ;
\draw [shift={(252.66,124.24)}, rotate = 108.95] [fill={rgb, 255:red, 0; green, 0; blue, 0 }  ][line width=0.08]  [draw opacity=0] (8.93,-4.29) -- (0,0) -- (8.93,4.29) -- cycle    ;
%Curve Lines [id:da7235025574921302] 
\draw    (239.2,167.15) .. controls (229.9,209.15) and (224.5,251.8) .. (254.41,269.96) .. controls (284.32,288.12) and (348.25,291.55) .. (440.2,293.2) ;
\draw [shift={(232.1,227.49)}, rotate = 267.82] [fill={rgb, 255:red, 0; green, 0; blue, 0 }  ][line width=0.08]  [draw opacity=0] (8.93,-4.29) -- (0,0) -- (8.93,4.29) -- cycle    ;
\draw [shift={(351.01,290.05)}, rotate = 183.96] [fill={rgb, 255:red, 0; green, 0; blue, 0 }  ][line width=0.08]  [draw opacity=0] (8.93,-4.29) -- (0,0) -- (8.93,4.29) -- cycle    ;
%Curve Lines [id:da6596914067421401] 
\draw    (118.2,90.15) .. controls (123.2,96.15) and (122.2,96.15) .. (126.2,103.15) ;
%Curve Lines [id:da23111158986903058] 
\draw    (132.2,115.15) .. controls (137.2,128.15) and (141.2,141.15) .. (143.2,157.15) ;
%Curve Lines [id:da9981033941356783] 
\draw    (144.2,172.15) .. controls (148.2,227.15) and (136.2,289.15) .. (120.27,309.12) ;
%Curve Lines [id:da23910891544892932] 
\draw  [dash pattern={on 4.5pt off 4.5pt}]  (253.2,77.05) .. controls (258.7,60.55) and (279.4,72.95) .. (284.2,91.55) .. controls (289,110.15) and (272.2,143.05) .. (268.7,154.05) .. controls (265.2,165.05) and (260.7,181.05) .. (255.7,191.55) .. controls (250.7,202.05) and (221.2,196.05) .. (219.7,177.05) .. controls (218.2,158.05) and (223.2,148.55) .. (234.2,122.05) .. controls (245.2,95.55) and (234.7,120.55) .. (253.2,77.05) -- cycle ;
%Curve Lines [id:da16098984403322236] 
\draw    (265.06,464.04) .. controls (277.88,454.14) and (279.53,457.16) .. (285.2,463.4) ;
\draw [shift={(262.7,465.9)}, rotate = 321.11] [fill={rgb, 255:red, 0; green, 0; blue, 0 }  ][line width=0.08]  [draw opacity=0] (8.93,-4.29) -- (0,0) -- (8.93,4.29) -- cycle    ;
%Curve Lines [id:da324176041715751] 
\draw    (257.22,484.47) .. controls (269.44,474.65) and (267.65,476.7) .. (274.7,481.4) ;
\draw [shift={(254.7,486.5)}, rotate = 321.11] [fill={rgb, 255:red, 0; green, 0; blue, 0 }  ][line width=0.08]  [draw opacity=0] (8.93,-4.29) -- (0,0) -- (8.93,4.29) -- cycle    ;
%Curve Lines [id:da7349134069511734] 
\draw    (241.16,452.94) .. controls (249.27,434.92) and (235.82,430.55) .. (228.7,433.4) ;
\draw [shift={(239.7,455.9)}, rotate = 298.22] [fill={rgb, 255:red, 0; green, 0; blue, 0 }  ][line width=0.08]  [draw opacity=0] (8.93,-4.29) -- (0,0) -- (8.93,4.29) -- cycle    ;
%Curve Lines [id:da6608683284097078] 
\draw    (231.59,476.01) .. controls (236.65,458.41) and (231.35,453.92) .. (224.7,454.4) ;
\draw [shift={(230.7,478.9)}, rotate = 288] [fill={rgb, 255:red, 0; green, 0; blue, 0 }  ][line width=0.08]  [draw opacity=0] (8.93,-4.29) -- (0,0) -- (8.93,4.29) -- cycle    ;

% Text Node
\draw (37,16) node [anchor=north west][inner sep=0.75pt]   [align=left] {修改前};
% Text Node
\draw (39,352) node [anchor=north west][inner sep=0.75pt]   [align=left] {修改后};
% Text Node
\draw (58,257.95) node [anchor=north west][inner sep=0.75pt]    {$V_{0}$};
% Text Node
\draw (216,280.95) node [anchor=north west][inner sep=0.75pt]    {$W$};
% Text Node
\draw (519,279.95) node [anchor=north west][inner sep=0.75pt]    {$V_{1}$};
% Text Node
\draw (263.68,101.2) node [anchor=north west][inner sep=0.75pt]    {$p'$};
% Text Node
\draw (247.68,148.4) node [anchor=north west][inner sep=0.75pt]    {$p$};


\end{tikzpicture}
}}
\newcommand{\FigFiveFive}{%
  \resizebox{!}{.44\textheight}{% Pattern Info
 
\tikzset{
pattern size/.store in=\mcSize, 
pattern size = 5pt,
pattern thickness/.store in=\mcThickness, 
pattern thickness = 0.3pt,
pattern radius/.store in=\mcRadius, 
pattern radius = 1pt}
\makeatletter
\pgfutil@ifundefined{pgf@pattern@name@_1qw8iehvq}{
\pgfdeclarepatternformonly[\mcThickness,\mcSize]{_1qw8iehvq}
{\pgfqpoint{0pt}{0pt}}
{\pgfpoint{\mcSize}{\mcSize}}
{\pgfpoint{\mcSize}{\mcSize}}
{
\pgfsetcolor{\tikz@pattern@color}
\pgfsetlinewidth{\mcThickness}
\pgfpathmoveto{\pgfqpoint{0pt}{\mcSize}}
\pgfpathlineto{\pgfpoint{\mcSize+\mcThickness}{-\mcThickness}}
\pgfpathmoveto{\pgfqpoint{0pt}{0pt}}
\pgfpathlineto{\pgfpoint{\mcSize+\mcThickness}{\mcSize+\mcThickness}}
\pgfusepath{stroke}
}}
\makeatother
\tikzset{every picture/.style={line width=0.75pt}} %set default line width to 0.75pt        

\begin{tikzpicture}[x=0.75pt,y=0.75pt,yscale=-1,xscale=1]
%uncomment if require: \path (0,724); %set diagram left start at 0, and has height of 724

%Shape: Arc [id:dp11600535545844648] 
\draw  [draw opacity=0] (120.27,309.12) .. controls (101.28,307.87) and (86,257.86) .. (86,196.35) .. controls (86,134.05) and (101.67,83.55) .. (121,83.55) -- (121,196.35) -- cycle ; \draw   (120.27,309.12) .. controls (101.28,307.87) and (86,257.86) .. (86,196.35) .. controls (86,134.05) and (101.67,83.55) .. (121,83.55) ;  
%Curve Lines [id:da8485262314692139] 
\draw    (121,83.55) .. controls (161.2,76.15) and (238.6,73.27) .. (264.4,72.67) .. controls (290.2,72.07) and (345.73,70.67) .. (356.4,71.33) .. controls (367.07,72) and (372.4,71.33) .. (381.73,78.67) .. controls (391.07,86) and (395.73,93.33) .. (397.73,106.67) .. controls (399.73,120) and (397.07,144.67) .. (381.07,148.67) .. controls (365.07,152.67) and (311.48,159) .. (280.4,165.33) .. controls (249.32,171.67) and (241.07,180) .. (240.4,190.67) .. controls (239.73,201.33) and (243.07,210.67) .. (255.73,214) .. controls (268.4,217.33) and (372.6,229.4) .. (414.6,234.4) ;
%Shape: Ellipse [id:dp27803344135238306] 
\draw   (390.5,276.7) .. controls (390.5,253.34) and (401.29,234.4) .. (414.6,234.4) .. controls (427.91,234.4) and (438.7,253.34) .. (438.7,276.7) .. controls (438.7,300.06) and (427.91,319) .. (414.6,319) .. controls (401.29,319) and (390.5,300.06) .. (390.5,276.7) -- cycle ;
%Curve Lines [id:da3705907763733033] 
\draw    (120.27,309.12) .. controls (160.2,317.15) and (322.6,337.4) .. (414.6,319) ;
%Curve Lines [id:da9981033941356783] 
\draw  [dash pattern={on 4.5pt off 4.5pt}]  (121,83.55) .. controls (153.73,114.67) and (145.73,308) .. (120.27,309.12) ;
%Curve Lines [id:da012807667955076152] 
\draw    (320.4,71.33) .. controls (315.07,91.33) and (324.4,147.33) .. (338.4,155.33) ;
%Curve Lines [id:da2813386689403522] 
\draw    (264.4,72.67) .. controls (260.27,89.8) and (260.4,147.33) .. (280.4,165.33) ;
%Shape: Circle [id:dp04851694771864945] 
\draw  [color={rgb, 255:red, 0; green, 0; blue, 0 }  ,draw opacity=1 ][fill={rgb, 255:red, 0; green, 0; blue, 0 }  ,fill opacity=1 ] (365.07,110.93) .. controls (365.07,109.98) and (365.84,109.2) .. (366.8,109.2) .. controls (367.76,109.2) and (368.53,109.98) .. (368.53,110.93) .. controls (368.53,111.89) and (367.76,112.67) .. (366.8,112.67) .. controls (365.84,112.67) and (365.07,111.89) .. (365.07,110.93) -- cycle ;
%Curve Lines [id:da2435713626297017] 
\draw    (86.53,199.87) .. controls (129.73,194) and (215.07,183.33) .. (240.4,190.67) ;
\draw [shift={(168.37,190.47)}, rotate = 174.98] [fill={rgb, 255:red, 0; green, 0; blue, 0 }  ][line width=0.08]  [draw opacity=0] (8.93,-4.29) -- (0,0) -- (8.93,4.29) -- cycle    ;
%Shape: Circle [id:dp10568652884971608] 
\draw  [color={rgb, 255:red, 0; green, 0; blue, 0 }  ,draw opacity=1 ][fill={rgb, 255:red, 0; green, 0; blue, 0 }  ,fill opacity=1 ] (213.73,187.6) .. controls (213.73,186.64) and (214.51,185.87) .. (215.47,185.87) .. controls (216.42,185.87) and (217.2,186.64) .. (217.2,187.6) .. controls (217.2,188.56) and (216.42,189.33) .. (215.47,189.33) .. controls (214.51,189.33) and (213.73,188.56) .. (213.73,187.6) -- cycle ;
%Curve Lines [id:da5244903731780326] 
\draw    (215.47,187.6) .. controls (226.67,153.33) and (292.73,136.6) .. (317.07,129.93) .. controls (341.4,123.27) and (350.07,117.6) .. (366.8,110.93) ;
\draw [shift={(264.5,147.05)}, rotate = 156.25] [fill={rgb, 255:red, 0; green, 0; blue, 0 }  ][line width=0.08]  [draw opacity=0] (8.93,-4.29) -- (0,0) -- (8.93,4.29) -- cycle    ;
\draw [shift={(346.95,119.51)}, rotate = 157.27] [fill={rgb, 255:red, 0; green, 0; blue, 0 }  ][line width=0.08]  [draw opacity=0] (8.93,-4.29) -- (0,0) -- (8.93,4.29) -- cycle    ;
%Shape: Circle [id:dp22311643944119108] 
\draw  [color={rgb, 255:red, 0; green, 0; blue, 0 }  ,draw opacity=1 ][fill={rgb, 255:red, 0; green, 0; blue, 0 }  ,fill opacity=1 ] (267.07,144.93) .. controls (267.07,143.98) and (267.84,143.2) .. (268.8,143.2) .. controls (269.76,143.2) and (270.53,143.98) .. (270.53,144.93) .. controls (270.53,145.89) and (269.76,146.67) .. (268.8,146.67) .. controls (267.84,146.67) and (267.07,145.89) .. (267.07,144.93) -- cycle ;
%Shape: Circle [id:dp04098981639747079] 
\draw  [color={rgb, 255:red, 0; green, 0; blue, 0 }  ,draw opacity=1 ][fill={rgb, 255:red, 0; green, 0; blue, 0 }  ,fill opacity=1 ] (322.4,127.93) .. controls (322.4,126.98) and (323.18,126.2) .. (324.13,126.2) .. controls (325.09,126.2) and (325.87,126.98) .. (325.87,127.93) .. controls (325.87,128.89) and (325.09,129.67) .. (324.13,129.67) .. controls (323.18,129.67) and (322.4,128.89) .. (322.4,127.93) -- cycle ;
%Curve Lines [id:da29083691744655693] 
\draw    (309.4,178.93) .. controls (305.97,169.58) and (300.5,153.53) .. (295.45,140.04) ;
\draw [shift={(294.4,137.27)}, rotate = 69.15] [fill={rgb, 255:red, 0; green, 0; blue, 0 }  ][line width=0.08]  [draw opacity=0] (10.72,-5.15) -- (0,0) -- (10.72,5.15) -- (7.12,0) -- cycle    ;
%Curve Lines [id:da731849974099766] 
\draw [color={rgb, 255:red, 128; green, 128; blue, 128 }  ,draw opacity=1 ]   (321.2,109.93) .. controls (330.07,107.27) and (346.4,102.93) .. (351.73,99.93) .. controls (357.07,96.93) and (366.4,92.6) .. (373.73,94.93) .. controls (381.07,97.27) and (383.73,104.6) .. (383.73,112.93) .. controls (383.73,121.27) and (379.07,129.93) .. (367.4,130.93) .. controls (355.73,131.93) and (344.07,134.27) .. (328.73,140.93) ;
%Curve Lines [id:da917710178586405] 
\draw [color={rgb, 255:red, 128; green, 128; blue, 128 }  ,draw opacity=1 ]   (196.2,176.27) .. controls (218.53,184.6) and (223.2,152.93) .. (263.2,122.93) ;
%Curve Lines [id:da9073880894350372] 
\draw [color={rgb, 255:red, 128; green, 128; blue, 128 }  ,draw opacity=1 ]   (196.2,176.27) .. controls (195.53,191.27) and (196.2,192.93) .. (195.53,203.27) ;
%Curve Lines [id:da5399692162948594] 
\draw [color={rgb, 255:red, 128; green, 128; blue, 128 }  ,draw opacity=1 ]   (195.53,203.27) .. controls (202.2,203.27) and (206.87,210.93) .. (205.53,218.6) ;
%Curve Lines [id:da47464187974425165] 
\draw [color={rgb, 255:red, 128; green, 128; blue, 128 }  ,draw opacity=1 ]   (205.53,218.6) .. controls (215.53,217.93) and (212.2,218.27) .. (224.87,218.27) ;
%Curve Lines [id:da32938792477673096] 
\draw [color={rgb, 255:red, 128; green, 128; blue, 128 }  ,draw opacity=1 ]   (234.53,202.93) .. controls (226.2,202.93) and (226.2,210.6) .. (224.87,218.27) ;
%Curve Lines [id:da3913110489912415] 
\draw [color={rgb, 255:red, 128; green, 128; blue, 128 }  ,draw opacity=1 ]   (234.53,202.93) .. controls (235.2,197.6) and (236.2,187.6) .. (236.87,182.6) ;
%Curve Lines [id:da25485336033856554] 
\draw [color={rgb, 255:red, 128; green, 128; blue, 128 }  ,draw opacity=1 ]   (236.87,182.6) .. controls (220.2,175.6) and (248.53,164.6) .. (272.2,153.6) ;
%Curve Lines [id:da6458925522829497] 
\draw    (272.1,324.9) .. controls (267.1,312.9) and (260.1,240.9) .. (283.9,218) ;
%Curve Lines [id:da961485426052637] 
\draw    (325.2,326) .. controls (320.2,314) and (314.6,247.9) .. (338.4,225) ;
%Straight Lines [id:da01763230108924163] 
\draw    (54.2,498.8) -- (210.2,498.8) -- (407.2,498.8) -- (556.2,498.8) ;
\draw [shift={(558.2,498.8)}, rotate = 180] [color={rgb, 255:red, 0; green, 0; blue, 0 }  ][line width=0.75]    (10.93,-3.29) .. controls (6.95,-1.4) and (3.31,-0.3) .. (0,0) .. controls (3.31,0.3) and (6.95,1.4) .. (10.93,3.29)   ;
\draw [shift={(137.2,498.8)}, rotate = 180] [fill={rgb, 255:red, 0; green, 0; blue, 0 }  ][line width=0.08]  [draw opacity=0] (8.93,-4.29) -- (0,0) -- (8.93,4.29) -- cycle    ;
\draw [shift={(313.7,498.8)}, rotate = 180] [fill={rgb, 255:red, 0; green, 0; blue, 0 }  ][line width=0.08]  [draw opacity=0] (8.93,-4.29) -- (0,0) -- (8.93,4.29) -- cycle    ;
\draw [shift={(487.7,498.8)}, rotate = 180] [fill={rgb, 255:red, 0; green, 0; blue, 0 }  ][line width=0.08]  [draw opacity=0] (8.93,-4.29) -- (0,0) -- (8.93,4.29) -- cycle    ;
%Straight Lines [id:da9914116345476351] 
\draw    (210.2,410.8) -- (210.2,586.8) ;
%Straight Lines [id:da6001205713203279] 
\draw    (407.2,410.8) -- (407.2,586.8) ;
%Straight Lines [id:da09988333183730846] 
\draw [color={rgb, 255:red, 128; green, 128; blue, 128 }  ,draw opacity=1 ]   (186.2,444.2) -- (233,444.2) ;
%Straight Lines [id:da029517889875373515] 
\draw [color={rgb, 255:red, 128; green, 128; blue, 128 }  ,draw opacity=1 ]   (186.2,555.2) -- (233,555.2) ;
%Straight Lines [id:da8135836549217623] 
\draw [color={rgb, 255:red, 128; green, 128; blue, 128 }  ,draw opacity=1 ]   (156.2,474.2) -- (156.2,521.8) ;
%Straight Lines [id:da7749898752951888] 
\draw [color={rgb, 255:red, 128; green, 128; blue, 128 }  ,draw opacity=1 ]   (265.2,474.8) -- (265.2,522.4) ;
%Curve Lines [id:da3250418176678912] 
\draw [color={rgb, 255:red, 128; green, 128; blue, 128 }  ,draw opacity=1 ]   (156.2,474.2) .. controls (185.2,474.8) and (186.2,473.8) .. (186.2,444.2) ;
%Curve Lines [id:da7866300470342965] 
\draw [color={rgb, 255:red, 128; green, 128; blue, 128 }  ,draw opacity=1 ]   (265.2,474.8) .. controls (233.2,474.8) and (233.2,470.8) .. (233,444.2) ;
%Curve Lines [id:da07248892086971026] 
\draw [color={rgb, 255:red, 128; green, 128; blue, 128 }  ,draw opacity=1 ]   (265.2,522.4) .. controls (233.2,524) and (232.2,520.8) .. (233,555.2) ;
%Curve Lines [id:da3985269331105422] 
\draw [color={rgb, 255:red, 128; green, 128; blue, 128 }  ,draw opacity=1 ]   (156.2,521.8) .. controls (185.2,522.4) and (187.2,521.8) .. (186.2,555.2) ;

%Shape: Rectangle [id:dp9149170955306568] 
\draw  [color={rgb, 255:red, 155; green, 155; blue, 155 }  ,draw opacity=1 ][pattern=_1qw8iehvq,pattern size=6pt,pattern thickness=0.75pt,pattern radius=0pt, pattern color={rgb, 255:red, 155; green, 155; blue, 155}][dash pattern={on 0.84pt off 2.51pt}] (265.2,487.8) -- (351.2,487.8) -- (351.2,509.4) -- (265.2,509.4) -- cycle ;
%Straight Lines [id:da03635805043286633] 
\draw    (265.2,487.8) -- (351.2,487.8) ;
\draw [shift={(315.2,487.8)}, rotate = 180] [fill={rgb, 255:red, 0; green, 0; blue, 0 }  ][line width=0.08]  [draw opacity=0] (12,-3) -- (0,0) -- (12,3) -- cycle    ;
%Straight Lines [id:da19547677615858094] 
\draw    (265.2,509.4) -- (351.2,509.4) ;
\draw [shift={(315.2,509.4)}, rotate = 180] [fill={rgb, 255:red, 0; green, 0; blue, 0 }  ][line width=0.08]  [draw opacity=0] (12,-3) -- (0,0) -- (12,3) -- cycle    ;
%Straight Lines [id:da8704658441545563] 
\draw [line width=2.25]    (265.2,487.8) -- (265.2,509.4) ;
%Straight Lines [id:da31683373404892223] 
\draw [line width=2.25]    (351.2,487.8) -- (351.2,509.4) ;
%Shape: Circle [id:dp9276595972959283] 
\draw  [color={rgb, 255:red, 0; green, 0; blue, 0 }  ,draw opacity=1 ][fill={rgb, 255:red, 0; green, 0; blue, 0 }  ,fill opacity=1 ] (208.47,498.8) .. controls (208.47,497.84) and (209.24,497.07) .. (210.2,497.07) .. controls (211.16,497.07) and (211.93,497.84) .. (211.93,498.8) .. controls (211.93,499.76) and (211.16,500.53) .. (210.2,500.53) .. controls (209.24,500.53) and (208.47,499.76) .. (208.47,498.8) -- cycle ;
%Shape: Circle [id:dp7467395915735052] 
\draw  [color={rgb, 255:red, 0; green, 0; blue, 0 }  ,draw opacity=1 ][fill={rgb, 255:red, 0; green, 0; blue, 0 }  ,fill opacity=1 ] (405.47,498.8) .. controls (405.47,497.84) and (406.24,497.07) .. (407.2,497.07) .. controls (408.16,497.07) and (408.93,497.84) .. (408.93,498.8) .. controls (408.93,499.76) and (408.16,500.53) .. (407.2,500.53) .. controls (406.24,500.53) and (405.47,499.76) .. (405.47,498.8) -- cycle ;
%Curve Lines [id:da265944666977898] 
\draw [color={rgb, 255:red, 128; green, 128; blue, 128 }  ,draw opacity=1 ]   (351.2,487.8) .. controls (351,480.4) and (351,489.9) .. (349,473.4) .. controls (347,456.9) and (363.5,471.4) .. (388.5,446.4) .. controls (413.5,421.4) and (439.5,448.4) .. (449.5,458.9) .. controls (459.5,469.4) and (460,524.4) .. (451,536.4) .. controls (442,548.4) and (412,575.4) .. (384,548.4) .. controls (356,521.4) and (349.5,538.4) .. (350.5,522.9) .. controls (351.5,507.4) and (350.5,518.4) .. (351.2,509.4) ;
%Curve Lines [id:da0031127069057219625] 
\draw    (220,459.4) .. controls (218.32,473.8) and (233.33,487.06) .. (247.55,487.91) ;
\draw [shift={(250.5,487.9)}, rotate = 176.19] [fill={rgb, 255:red, 0; green, 0; blue, 0 }  ][line width=0.08]  [draw opacity=0] (8.93,-4.29) -- (0,0) -- (8.93,4.29) -- cycle    ;
%Curve Lines [id:da7530784821797115] 
\draw    (217.5,533.4) .. controls (217.03,515.54) and (233.79,509.17) .. (249.51,509.71) ;
\draw [shift={(252.5,509.9)}, rotate = 185.19] [fill={rgb, 255:red, 0; green, 0; blue, 0 }  ][line width=0.08]  [draw opacity=0] (8.93,-4.29) -- (0,0) -- (8.93,4.29) -- cycle    ;
%Straight Lines [id:da6274095608076921] 
\draw    (368.5,474.9) -- (392.22,488.97) ;
\draw [shift={(394.8,490.5)}, rotate = 210.67] [fill={rgb, 255:red, 0; green, 0; blue, 0 }  ][line width=0.08]  [draw opacity=0] (8.93,-4.29) -- (0,0) -- (8.93,4.29) -- cycle    ;
%Straight Lines [id:da3195351705698153] 
\draw    (368,522.4) -- (391.22,508.54) ;
\draw [shift={(393.8,507)}, rotate = 149.17] [fill={rgb, 255:red, 0; green, 0; blue, 0 }  ][line width=0.08]  [draw opacity=0] (8.93,-4.29) -- (0,0) -- (8.93,4.29) -- cycle    ;
%Straight Lines [id:da3173952133692499] 
\draw    (339.5,541.4) -- (346.66,516.78) ;
\draw [shift={(347.5,513.9)}, rotate = 106.22] [fill={rgb, 255:red, 0; green, 0; blue, 0 }  ][line width=0.08]  [draw opacity=0] (8.93,-4.29) -- (0,0) -- (8.93,4.29) -- cycle    ;
%Straight Lines [id:da8068004311628002] 
\draw    (277,541.4) -- (270.3,517.29) ;
\draw [shift={(269.5,514.4)}, rotate = 74.48] [fill={rgb, 255:red, 0; green, 0; blue, 0 }  ][line width=0.08]  [draw opacity=0] (8.93,-4.29) -- (0,0) -- (8.93,4.29) -- cycle    ;
%Curve Lines [id:da6363495959445716] 
\draw    (161.2,446.2) .. controls (150.31,428.38) and (182.54,278.25) .. (199.49,214.5) ;
\draw [shift={(200,212.6)}, rotate = 105.02] [color={rgb, 255:red, 0; green, 0; blue, 0 }  ][line width=0.75]    (10.93,-3.29) .. controls (6.95,-1.4) and (3.31,-0.3) .. (0,0) .. controls (3.31,0.3) and (6.95,1.4) .. (10.93,3.29)   ;
%Curve Lines [id:da629840018543867] 
\draw    (456.2,443.2) .. controls (477.09,417.33) and (467.3,154.84) .. (386.42,121.68) ;
\draw [shift={(385.2,121.2)}, rotate = 20.71] [color={rgb, 255:red, 0; green, 0; blue, 0 }  ][line width=0.75]    (10.93,-3.29) .. controls (6.95,-1.4) and (3.31,-0.3) .. (0,0) .. controls (3.31,0.3) and (6.95,1.4) .. (10.93,3.29)   ;

% Text Node
\draw (209.87,187) node [anchor=north west][inner sep=0.75pt]    {$p$};
% Text Node
\draw (268.87,117.67) node [anchor=north west][inner sep=0.75pt]    {$p_{1}$};
% Text Node
\draw (305.53,109.33) node [anchor=north west][inner sep=0.75pt]    {$p_{2}$};
% Text Node
\draw (363.2,109) node [anchor=north west][inner sep=0.75pt]    {$p'$};
% Text Node
\draw (309.2,180.33) node [anchor=north west][inner sep=0.75pt]    {$T$};
% Text Node
\draw (244.4,333.4) node [anchor=north west][inner sep=0.75pt]    {$f^{-1}( b_{1})$};
% Text Node
\draw (310.9,332.8) node [anchor=north west][inner sep=0.75pt]    {$f^{-1}( b_{2})$};
% Text Node
\draw (195.8,504.9) node [anchor=north west][inner sep=0.75pt]    {$0$};
% Text Node
\draw (410.93,502.2) node [anchor=north west][inner sep=0.75pt]    {$e$};
% Text Node
\draw (268.8,545.4) node [anchor=north west][inner sep=0.75pt]    {$U_{0}$};
% Text Node
\draw (330.8,545.4) node [anchor=north west][inner sep=0.75pt]    {$U_{1}$};
% Text Node
\draw (299.3,457.4) node [anchor=north west][inner sep=0.75pt]    {$L_{0}$};
% Text Node
\draw (181.3,475.9) node [anchor=north west][inner sep=0.75pt]    {$L_{1}$};
% Text Node
\draw (425.3,472.4) node [anchor=north west][inner sep=0.75pt]    {$L_{2}$};
% Text Node
\draw (542,473.4) node [anchor=north west][inner sep=0.75pt]    {$x_{1}$};
% Text Node
\draw (141,356) node [anchor=north west][inner sep=0.75pt]    {$g_{1}$};
% Text Node
\draw (475,331) node [anchor=north west][inner sep=0.75pt]    {$g_{2}$};


\end{tikzpicture}
}}
\newcommand{\FigFiveThree}{%
\begin{tikzpicture}[>=Stealth,x=6cm,y=3.2cm,line cap=round,line join=round]
  \draw[->] (-.28,0)--(1.29,0) node[right] {$x_1$};
  \draw[->] (0,-.29)--(0,.57) node[above] {$y$};

  % 在 0、1 附近分别沿 y=x、y=1-x：两处零点准确，且 |v'|=1。
  \draw[thick]
    (-.24,-.24)--(.12,.12)
    .. controls (.22,.22) and (.36,.39) .. (.50,.39)
    .. controls (.64,.39) and (.78,.22) .. (.88,.12)
    --(1.24,-.24);
  \node[above right=1pt] at (.53,.39) {$y=v(x_1)$};
  \node[below right=2pt] at (0,0) {$0$};
  \node[below left=2pt] at (1,0) {$1$};
\end{tikzpicture}}

\newcommand{\FigFiveFour}{%
\begin{tikzpicture}[>=Stealth,x=6cm,y=3.2cm,line cap=round,line join=round]
  \draw[->] (-.28,0)--(1.29,0) node[right] {$x_1$};
  \draw[->] (0,-.35)--(0,.57) node[above] {$y$};

  % 原函数：在 (0,1) 内为正，在 0、1 处为零，区间外为负。
  \draw[thick]
    (-.24,-.24)--(.12,.12)
    .. controls (.22,.22) and (.36,.39) .. (.50,.39)
    .. controls (.64,.39) and (.78,.22) .. (.88,.12)
    --(1.24,-.24);
  \node[above right=1pt] at (.53,.39) {$y=v(x_1)$};

  % 修改后的轴上函数处处为负；在中央紧邻域之外与 v 重合。
  \draw[thick]
    (-.24,-.24)--(-.18,-.18)
    .. controls (-.08,-.08) and (1.08,-.08) .. (1.18,-.18)
    --(1.24,-.24);
  \node[below=3pt] at (.50,-.105) {$y=v'(x_1,0)$};
  \node[below right=2pt] at (0,0) {$0$};
  \node[below left=2pt] at (1,0) {$1$};
\end{tikzpicture}}

\newcommand{\FigFiveSix}{%
  \resizebox{.75\textwidth}{!}{\tikzset{every picture/.style={line width=0.75pt}} %set default line width to 0.75pt        

\begin{tikzpicture}[x=0.75pt,y=0.75pt,yscale=-1,xscale=1]
%uncomment if require: \path (0,724); %set diagram left start at 0, and has height of 724

%Straight Lines [id:da3129485065184029] 
\draw    (123,254.6) -- (551.2,254.6) ;
\draw [shift={(553.2,254.6)}, rotate = 180] [color={rgb, 255:red, 0; green, 0; blue, 0 }  ][line width=0.75]    (10.93,-3.29) .. controls (6.95,-1.4) and (3.31,-0.3) .. (0,0) .. controls (3.31,0.3) and (6.95,1.4) .. (10.93,3.29)   ;
%Straight Lines [id:da40717068394496947] 
\draw    (338,406.6) -- (338,91.2) ;
\draw [shift={(338,89.2)}, rotate = 90] [color={rgb, 255:red, 0; green, 0; blue, 0 }  ][line width=0.75]    (10.93,-3.29) .. controls (6.95,-1.4) and (3.31,-0.3) .. (0,0) .. controls (3.31,0.3) and (6.95,1.4) .. (10.93,3.29)   ;
%Shape: Polygon Curved [id:ds3311728239547659] 
\draw  [fill={rgb, 255:red, 155; green, 155; blue, 155 }  ,fill opacity=0.32 ] (258.2,161.2) .. controls (276.2,143.2) and (359.2,116.2) .. (390.2,136.2) .. controls (421.2,156.2) and (451.2,167.2) .. (447.2,187.2) .. controls (443.2,207.2) and (445.2,249.2) .. (447.2,269.2) .. controls (449.2,289.2) and (426.7,351.7) .. (399.2,363.2) .. controls (371.7,374.7) and (320.82,359.58) .. (295.2,353.2) .. controls (269.57,346.83) and (200.2,343.2) .. (209.2,314.2) .. controls (218.2,285.2) and (234.2,273.2) .. (230.2,237.2) .. controls (226.2,201.2) and (240.2,179.2) .. (258.2,161.2) -- cycle ;
%Curve Lines [id:da06371900121051755] 
\draw    (168.2,395.2) .. controls (219.2,365.2) and (280.2,329.2) .. (299.2,308.2) .. controls (318.2,287.2) and (328.2,275.2) .. (338.1,254.6) .. controls (348,234) and (368.2,215.2) .. (389.2,199.2) .. controls (410.2,183.2) and (500.2,152.2) .. (538.2,139.2) ;
%Curve Lines [id:da6348921142501356] 
\draw    (168.2,395.2) .. controls (209.2,370.2) and (220.2,365.2) .. (260.2,339.2) .. controls (300.2,313.2) and (291.2,306.2) .. (294.2,294.2) .. controls (297.2,282.2) and (297,254.4) .. (311.5,254.4) .. controls (326,254.4) and (332.4,254.4) .. (338.1,254.6) .. controls (343.8,254.8) and (344.4,254.4) .. (363.9,254.9) .. controls (383.4,255.4) and (387.5,242.4) .. (398,213.4) .. controls (408.5,184.4) and (438.4,174.9) .. (465.9,164.4) .. controls (493.4,153.9) and (507.9,149.4) .. (538.2,139.2)(174.88,386.47) -- (179.01,393.32)(183.27,381.43) -- (187.37,388.3)(192.21,376.12) -- (196.28,383)(200.73,371.09) -- (204.78,377.99)(209.45,365.96) -- (213.51,372.85)(217.94,360.92) -- (222.05,367.78)(226.28,355.89) -- (230.45,362.72)(234.82,350.63) -- (239.05,357.42)(243.3,345.3) -- (247.58,352.06)(251.69,339.94) -- (256.02,346.67)(260.33,334.33) -- (264.75,341)(268.44,328.71) -- (273.15,335.18)(275.89,322.92) -- (281.04,329.04)(282.59,316.59) -- (288.5,321.99)(287.62,309.4) -- (294.89,312.73)(289.34,301.08) -- (297.33,301.56)(290.86,290.72) -- (298.72,292.2)(292.36,281.03) -- (300.28,282.21)(294.21,270.71) -- (302.01,272.47)(297.6,260.38) -- (304.8,263.87)(306.23,251.42) -- (309.27,258.81)(317.69,250.4) -- (317.68,258.4)(327.75,250.42) -- (327.7,258.42)(337.8,250.59) -- (337.54,258.58)(347.68,250.64) -- (347.67,258.64)(357.73,250.76) -- (357.58,258.76)(367.46,250.82) -- (368.1,258.79)(375.29,248.54) -- (379.36,255.42)(381.05,242.92) -- (387.6,247.51)(385.53,234.75) -- (392.81,238.06)(389.22,225.75) -- (396.7,228.59)(392.71,216.27) -- (400.23,218.99)(396.48,206.66) -- (403.71,210.11)(401.95,197.5) -- (408.41,202.21)(409.03,189.42) -- (414.57,195.18)(417.21,182.65) -- (421.87,189.16)(425.83,177.19) -- (429.77,184.15)(435.39,172.31) -- (438.76,179.56)(444.47,168.37) -- (447.5,175.77)(454.09,164.59) -- (456.94,172.07)(463.24,161.13) -- (466.08,168.61)(472.77,157.53) -- (475.57,165.02)(482.14,154.08) -- (484.87,161.6)(491.55,150.73) -- (494.2,158.28)(501.01,147.46) -- (503.6,155.03)(510.37,144.29) -- (512.91,151.88)(520.04,141.06) -- (522.58,148.65)(529.53,137.89) -- (532.07,145.48) ;
%Curve Lines [id:da3876903888663663] 
\draw [color={rgb, 255:red, 155; green, 155; blue, 155 }  ,draw opacity=1 ]   (201.2,205.6) .. controls (202.36,240.13) and (263.37,280.12) .. (284.37,284.31) ;
\draw [shift={(286.2,284.6)}, rotate = 186.01] [color={rgb, 255:red, 155; green, 155; blue, 155 }  ,draw opacity=1 ][line width=0.75]    (10.93,-3.29) .. controls (6.95,-1.4) and (3.31,-0.3) .. (0,0) .. controls (3.31,0.3) and (6.95,1.4) .. (10.93,3.29)   ;
%Curve Lines [id:da04910207175146397] 
\draw [color={rgb, 255:red, 155; green, 155; blue, 155 }  ,draw opacity=1 ]   (417.2,118.6) .. controls (391.98,144.79) and (362.06,185.09) .. (367.58,203) ;
\draw [shift={(368.2,204.6)}, rotate = 244.8] [color={rgb, 255:red, 155; green, 155; blue, 155 }  ,draw opacity=1 ][line width=0.75]    (10.93,-3.29) .. controls (6.95,-1.4) and (3.31,-0.3) .. (0,0) .. controls (3.31,0.3) and (6.95,1.4) .. (10.93,3.29)   ;

% Text Node
\draw (559,272.4) node [anchor=north west][inner sep=0.75pt]    {$\mathbb{R}^{a}$};
% Text Node
\draw (298,77.4) node [anchor=north west][inner sep=0.75pt]    {$\mathbb{R}^{b}$};
% Text Node
\draw (415,87.4) node [anchor=north west][inner sep=0.75pt]    {$h\left(\mathbb{R}^{a}\right)$};
% Text Node
\draw (165,174.4) node [anchor=north west][inner sep=0.75pt]    {$h_{1}^{'}\left(\mathbb{R}^{b}\right)$};
% Text Node
\draw (315,224.4) node [anchor=north west][inner sep=0.75pt]    {$O$};
% Text Node
\draw (371,298.4) node [anchor=north west][inner sep=0.75pt]    {$h( N)$};


\end{tikzpicture}
}}
\newcommand{\FigFiveEight}{%
  \resizebox{.75\textwidth}{!}{\tikzset{every picture/.style={line width=0.75pt}} %set default line width to 0.75pt        

\begin{tikzpicture}[x=0.75pt,y=0.75pt,yscale=-1,xscale=1]
%uncomment if require: \path (0,724); %set diagram left start at 0, and has height of 724

%Shape: Rectangle [id:dp07077866533791832] 
\draw   (158,22.6) -- (462.2,22.6) -- (462.2,513.6) -- (158,513.6) -- cycle ;
%Straight Lines [id:da9851204634935387] 
\draw    (158,268.1) -- (462.2,268.1) ;
%Straight Lines [id:da43558910082487867] 
\draw    (158,94.2) -- (462.2,94.2) ;
%Straight Lines [id:da17448914126821513] 
\draw    (158,442) -- (462.2,442) ;

%Straight Lines [id:da47802169517293325] 
\draw    (499.2,23.6) -- (561.2,23.6) ;
%Straight Lines [id:da7457038291310528] 
\draw    (561.2,249.1) -- (561.2,23.6) ;
\draw [shift={(561.2,23.6)}, rotate = 90] [color={rgb, 255:red, 0; green, 0; blue, 0 }  ][line width=0.75]    (0,5.59) -- (0,-5.59)(10.93,-4.9) .. controls (6.95,-2.3) and (3.31,-0.67) .. (0,0) .. controls (3.31,0.67) and (6.95,2.3) .. (10.93,4.9)   ;
%Straight Lines [id:da650969957269567] 
\draw    (561.2,289.1) -- (561.2,514.6) ;
\draw [shift={(561.2,514.6)}, rotate = 270] [color={rgb, 255:red, 0; green, 0; blue, 0 }  ][line width=0.75]    (0,5.59) -- (0,-5.59)(10.93,-4.9) .. controls (6.95,-2.3) and (3.31,-0.67) .. (0,0) .. controls (3.31,0.67) and (6.95,2.3) .. (10.93,4.9)   ;
%Straight Lines [id:da7935270356364614] 
\draw    (499.2,514.6) -- (561.2,514.6) ;
%Straight Lines [id:da47386519470126787] 
\draw    (74,95.2) -- (136,95.2) ;
%Straight Lines [id:da8984106025441775] 
\draw    (74,443) -- (136,443) ;
%Straight Lines [id:da16724332871276604] 
\draw    (74,248.6) -- (74,95.2) ;
\draw [shift={(74,95.2)}, rotate = 90] [color={rgb, 255:red, 0; green, 0; blue, 0 }  ][line width=0.75]    (0,5.59) -- (0,-5.59)(10.93,-4.9) .. controls (6.95,-2.3) and (3.31,-0.67) .. (0,0) .. controls (3.31,0.67) and (6.95,2.3) .. (10.93,4.9)   ;
%Straight Lines [id:da822207148295593] 
\draw    (74,443) -- (74,289.6) ;
\draw [shift={(74,443)}, rotate = 270] [color={rgb, 255:red, 0; green, 0; blue, 0 }  ][line width=0.75]    (0,5.59) -- (0,-5.59)(10.93,-4.9) .. controls (6.95,-2.3) and (3.31,-0.67) .. (0,0) .. controls (3.31,0.67) and (6.95,2.3) .. (10.93,4.9)   ;
%Curve Lines [id:da5273318006635597] 
\draw    (158.2,144.6) .. controls (198.4,143.2) and (220.2,143.6) .. (257.2,122.6) .. controls (294.2,101.6) and (330.2,156.6) .. (342.2,169.6) .. controls (354.2,182.6) and (387.2,208.6) .. (403.2,214.6) .. controls (419.2,220.6) and (445.2,233.6) .. (462.2,232.6) ;
%Curve Lines [id:da4648166730602935] 
\draw    (158.2,335.6) .. controls (198.4,334.2) and (224.2,328.6) .. (261.2,354.6) .. controls (298.2,380.6) and (327.2,384.6) .. (343.2,383.6) .. controls (359.2,382.6) and (381.2,371.6) .. (401.2,364.6) .. controls (421.2,357.6) and (443.2,357.6) .. (462.2,357.6) ;
%Straight Lines [id:da6273220203407115] 
\draw    (184,42) -- (248.2,42) ;
\draw [shift={(251.2,42)}, rotate = 180] [fill={rgb, 255:red, 0; green, 0; blue, 0 }  ][line width=0.08]  [draw opacity=0] (8.93,-4.29) -- (0,0) -- (8.93,4.29) -- cycle    ;
%Straight Lines [id:da461175668219148] 
\draw    (277.4,42) -- (341.6,42) ;
\draw [shift={(344.6,42)}, rotate = 180] [fill={rgb, 255:red, 0; green, 0; blue, 0 }  ][line width=0.08]  [draw opacity=0] (8.93,-4.29) -- (0,0) -- (8.93,4.29) -- cycle    ;
%Straight Lines [id:da012078210586639826] 
\draw    (370.8,42) -- (435,42) ;
\draw [shift={(438,42)}, rotate = 180] [fill={rgb, 255:red, 0; green, 0; blue, 0 }  ][line width=0.08]  [draw opacity=0] (8.93,-4.29) -- (0,0) -- (8.93,4.29) -- cycle    ;
%Straight Lines [id:da26997671030191983] 
\draw    (184.4,74.6) -- (248.6,74.6) ;
\draw [shift={(251.6,74.6)}, rotate = 180] [fill={rgb, 255:red, 0; green, 0; blue, 0 }  ][line width=0.08]  [draw opacity=0] (8.93,-4.29) -- (0,0) -- (8.93,4.29) -- cycle    ;
%Straight Lines [id:da5773404430323953] 
\draw    (277.8,74.6) -- (342,74.6) ;
\draw [shift={(345,74.6)}, rotate = 180] [fill={rgb, 255:red, 0; green, 0; blue, 0 }  ][line width=0.08]  [draw opacity=0] (8.93,-4.29) -- (0,0) -- (8.93,4.29) -- cycle    ;
%Straight Lines [id:da09971016989561443] 
\draw    (371.2,74.6) -- (435.4,74.6) ;
\draw [shift={(438.4,74.6)}, rotate = 180] [fill={rgb, 255:red, 0; green, 0; blue, 0 }  ][line width=0.08]  [draw opacity=0] (8.93,-4.29) -- (0,0) -- (8.93,4.29) -- cycle    ;
%Straight Lines [id:da2404403571150111] 
\draw    (184,463) -- (248.2,463) ;
\draw [shift={(251.2,463)}, rotate = 180] [fill={rgb, 255:red, 0; green, 0; blue, 0 }  ][line width=0.08]  [draw opacity=0] (8.93,-4.29) -- (0,0) -- (8.93,4.29) -- cycle    ;
%Straight Lines [id:da9088580522294909] 
\draw    (277.4,463) -- (341.6,463) ;
\draw [shift={(344.6,463)}, rotate = 180] [fill={rgb, 255:red, 0; green, 0; blue, 0 }  ][line width=0.08]  [draw opacity=0] (8.93,-4.29) -- (0,0) -- (8.93,4.29) -- cycle    ;
%Straight Lines [id:da6946377409910418] 
\draw    (370.8,463) -- (435,463) ;
\draw [shift={(438,463)}, rotate = 180] [fill={rgb, 255:red, 0; green, 0; blue, 0 }  ][line width=0.08]  [draw opacity=0] (8.93,-4.29) -- (0,0) -- (8.93,4.29) -- cycle    ;
%Straight Lines [id:da7174818861527796] 
\draw    (184.4,495.6) -- (248.6,495.6) ;
\draw [shift={(251.6,495.6)}, rotate = 180] [fill={rgb, 255:red, 0; green, 0; blue, 0 }  ][line width=0.08]  [draw opacity=0] (8.93,-4.29) -- (0,0) -- (8.93,4.29) -- cycle    ;
%Straight Lines [id:da8717209985128664] 
\draw    (277.8,495.6) -- (342,495.6) ;
\draw [shift={(345,495.6)}, rotate = 180] [fill={rgb, 255:red, 0; green, 0; blue, 0 }  ][line width=0.08]  [draw opacity=0] (8.93,-4.29) -- (0,0) -- (8.93,4.29) -- cycle    ;
%Straight Lines [id:da3661502137551441] 
\draw    (371.2,495.6) -- (435.4,495.6) ;
\draw [shift={(438.4,495.6)}, rotate = 180] [fill={rgb, 255:red, 0; green, 0; blue, 0 }  ][line width=0.08]  [draw opacity=0] (8.93,-4.29) -- (0,0) -- (8.93,4.29) -- cycle    ;
%Straight Lines [id:da5561992367865579] 
\draw    (184,283) -- (248.2,283) ;
\draw [shift={(251.2,283)}, rotate = 180] [fill={rgb, 255:red, 0; green, 0; blue, 0 }  ][line width=0.08]  [draw opacity=0] (8.93,-4.29) -- (0,0) -- (8.93,4.29) -- cycle    ;
%Straight Lines [id:da37458214795769207] 
\draw    (277.4,283) -- (341.6,283) ;
\draw [shift={(344.6,283)}, rotate = 180] [fill={rgb, 255:red, 0; green, 0; blue, 0 }  ][line width=0.08]  [draw opacity=0] (8.93,-4.29) -- (0,0) -- (8.93,4.29) -- cycle    ;
%Straight Lines [id:da201064259578437] 
\draw    (370.8,283) -- (435,283) ;
\draw [shift={(438,283)}, rotate = 180] [fill={rgb, 255:red, 0; green, 0; blue, 0 }  ][line width=0.08]  [draw opacity=0] (8.93,-4.29) -- (0,0) -- (8.93,4.29) -- cycle    ;
%Straight Lines [id:da7028644476904002] 
\draw    (185,252) -- (249.2,252) ;
\draw [shift={(252.2,252)}, rotate = 180] [fill={rgb, 255:red, 0; green, 0; blue, 0 }  ][line width=0.08]  [draw opacity=0] (8.93,-4.29) -- (0,0) -- (8.93,4.29) -- cycle    ;
%Straight Lines [id:da2819070251296256] 
\draw    (278.4,252) -- (342.6,252) ;
\draw [shift={(345.6,252)}, rotate = 180] [fill={rgb, 255:red, 0; green, 0; blue, 0 }  ][line width=0.08]  [draw opacity=0] (8.93,-4.29) -- (0,0) -- (8.93,4.29) -- cycle    ;
%Straight Lines [id:da2478802459231405] 
\draw    (371.8,252) -- (436,252) ;
\draw [shift={(439,252)}, rotate = 180] [fill={rgb, 255:red, 0; green, 0; blue, 0 }  ][line width=0.08]  [draw opacity=0] (8.93,-4.29) -- (0,0) -- (8.93,4.29) -- cycle    ;
%Curve Lines [id:da7196205574205272] 
\draw    (185,167) .. controls (219.93,167.58) and (223.02,162.01) .. (246.56,152.64) ;
\draw [shift={(249.2,151.6)}, rotate = 158.96] [fill={rgb, 255:red, 0; green, 0; blue, 0 }  ][line width=0.08]  [draw opacity=0] (8.93,-4.29) -- (0,0) -- (8.93,4.29) -- cycle    ;
%Curve Lines [id:da05418137699092462] 
\draw    (273.2,143.6) .. controls (302.68,127.16) and (324.67,173.94) .. (336.28,188.42) ;
\draw [shift={(338.2,190.6)}, rotate = 225] [fill={rgb, 255:red, 0; green, 0; blue, 0 }  ][line width=0.08]  [draw opacity=0] (8.93,-4.29) -- (0,0) -- (8.93,4.29) -- cycle    ;
%Curve Lines [id:da30736383462393846] 
\draw    (369,212) .. controls (384.39,226.82) and (410.43,236.59) .. (429.29,238.38) ;
\draw [shift={(432.2,238.6)}, rotate = 183.01] [fill={rgb, 255:red, 0; green, 0; blue, 0 }  ][line width=0.08]  [draw opacity=0] (8.93,-4.29) -- (0,0) -- (8.93,4.29) -- cycle    ;
%Curve Lines [id:da5221090932955107] 
\draw    (185.2,316.6) .. controls (212.6,313.76) and (235.54,320.75) .. (248.96,329.97) ;
\draw [shift={(251.2,331.6)}, rotate = 217.57] [fill={rgb, 255:red, 0; green, 0; blue, 0 }  ][line width=0.08]  [draw opacity=0] (8.93,-4.29) -- (0,0) -- (8.93,4.29) -- cycle    ;
%Curve Lines [id:da7989624514256152] 
\draw    (275,338) .. controls (290.39,352.82) and (316.43,362.59) .. (335.29,364.38) ;
\draw [shift={(338.2,364.6)}, rotate = 183.01] [fill={rgb, 255:red, 0; green, 0; blue, 0 }  ][line width=0.08]  [draw opacity=0] (8.93,-4.29) -- (0,0) -- (8.93,4.29) -- cycle    ;
%Curve Lines [id:da7695791716721757] 
\draw    (373.2,356.6) .. controls (395.16,347.05) and (415.31,337.5) .. (436.24,334.92) ;
\draw [shift={(439.2,334.6)}, rotate = 174.81] [fill={rgb, 255:red, 0; green, 0; blue, 0 }  ][line width=0.08]  [draw opacity=0] (8.93,-4.29) -- (0,0) -- (8.93,4.29) -- cycle    ;
%Straight Lines [id:da03983740989911888] 
\draw    (110.2,144.6) -- (136.2,144.6) ;
%Straight Lines [id:da05441593592501548] 
\draw    (136.2,144.6) -- (136.2,335.6) ;
%Straight Lines [id:da7547991935085584] 
\draw    (110.2,221.6) -- (110.2,144.6) ;
\draw [shift={(110.2,144.6)}, rotate = 90] [color={rgb, 255:red, 0; green, 0; blue, 0 }  ][line width=0.75]    (0,5.59) -- (0,-5.59)(10.93,-4.9) .. controls (6.95,-2.3) and (3.31,-0.67) .. (0,0) .. controls (3.31,0.67) and (6.95,2.3) .. (10.93,4.9)   ;
%Straight Lines [id:da9798694334860891] 
\draw    (110.2,335.6) -- (136.2,335.6) ;
%Straight Lines [id:da48287597712037555] 
\draw    (110.2,335.6) -- (110.2,258.6) ;
\draw [shift={(110.2,335.6)}, rotate = 270] [color={rgb, 255:red, 0; green, 0; blue, 0 }  ][line width=0.75]    (0,5.59) -- (0,-5.59)(10.93,-4.9) .. controls (6.95,-2.3) and (3.31,-0.67) .. (0,0) .. controls (3.31,0.67) and (6.95,2.3) .. (10.93,4.9)   ;

% Text Node
\draw (536,258.4) node [anchor=north west][inner sep=0.75pt]    {$h\left(\mathbb{R}^{n}\right)$};
% Text Node
\draw (59,262.4) node [anchor=north west][inner sep=0.75pt]    {$h( E)$};
% Text Node
\draw (91,228.4) node [anchor=north west][inner sep=0.75pt]    {$h(\overline{E}_{1})$};


\end{tikzpicture}
}}
\newcommand{\FigFiveNine}{%
  \resizebox{.75\textwidth}{!}{\tikzset{every picture/.style={line width=0.75pt}} %set default line width to 0.75pt        

\begin{tikzpicture}[x=0.75pt,y=0.75pt,yscale=-1,xscale=1]
%uncomment if require: \path (0,724); %set diagram left start at 0, and has height of 724

%Straight Lines [id:da3129485065184029] 
\draw    (123,254.6) -- (551.2,254.6) ;
\draw [shift={(553.2,254.6)}, rotate = 180] [color={rgb, 255:red, 0; green, 0; blue, 0 }  ][line width=0.75]    (10.93,-3.29) .. controls (6.95,-1.4) and (3.31,-0.3) .. (0,0) .. controls (3.31,0.3) and (6.95,1.4) .. (10.93,3.29)   ;
%Straight Lines [id:da40717068394496947] 
\draw    (338,406.6) -- (338,91.2) ;
\draw [shift={(338,89.2)}, rotate = 90] [color={rgb, 255:red, 0; green, 0; blue, 0 }  ][line width=0.75]    (10.93,-3.29) .. controls (6.95,-1.4) and (3.31,-0.3) .. (0,0) .. controls (3.31,0.3) and (6.95,1.4) .. (10.93,3.29)   ;
%Shape: Polygon Curved [id:ds3311728239547659] 
\draw  [fill={rgb, 255:red, 155; green, 155; blue, 155 }  ,fill opacity=0.32 ] (258.2,161.2) .. controls (276.2,143.2) and (359.2,116.2) .. (390.2,136.2) .. controls (421.2,156.2) and (451.2,167.2) .. (447.2,187.2) .. controls (443.2,207.2) and (445.2,249.2) .. (447.2,269.2) .. controls (449.2,289.2) and (426.7,351.7) .. (399.2,363.2) .. controls (371.7,374.7) and (320.82,359.58) .. (295.2,353.2) .. controls (269.57,346.83) and (200.2,343.2) .. (209.2,314.2) .. controls (218.2,285.2) and (234.2,273.2) .. (230.2,237.2) .. controls (226.2,201.2) and (240.2,179.2) .. (258.2,161.2) -- cycle ;
%Curve Lines [id:da06371900121051755] 
\draw    (168.2,395.2) .. controls (219.2,365.2) and (280.2,329.2) .. (299.2,308.2) .. controls (318.2,287.2) and (328.2,275.2) .. (338.1,254.6) .. controls (348,234) and (368.2,215.2) .. (389.2,199.2) .. controls (410.2,183.2) and (500.2,152.2) .. (538.2,139.2) ;
%Curve Lines [id:da6348921142501356] 
\draw    (168.2,395.2) .. controls (209.2,370.2) and (220.2,365.2) .. (260.2,339.2) .. controls (300.2,313.2) and (291.2,306.2) .. (294.2,294.2) .. controls (297.2,282.2) and (297,254.4) .. (311.5,254.4) .. controls (326,254.4) and (332.4,254.4) .. (338.1,254.6) .. controls (343.8,254.8) and (344.4,254.4) .. (363.9,254.9) .. controls (383.4,255.4) and (392.7,257.6) .. (403.2,228.6) .. controls (413.7,199.6) and (354.41,209.87) .. (324.2,200.6) .. controls (293.99,191.33) and (310.61,169.03) .. (320.2,163.6) .. controls (329.79,158.17) and (367.2,176.6) .. (402.2,177.6) .. controls (437.2,178.6) and (459.03,167.03) .. (465.9,164.4) .. controls (472.78,161.78) and (507.9,149.4) .. (538.2,139.2)(174.88,386.47) -- (179.01,393.32)(183.27,381.43) -- (187.37,388.3)(192.21,376.12) -- (196.28,383)(200.73,371.09) -- (204.78,377.99)(209.45,365.96) -- (213.51,372.85)(217.94,360.92) -- (222.05,367.78)(226.28,355.89) -- (230.45,362.72)(234.82,350.63) -- (239.05,357.42)(243.3,345.3) -- (247.58,352.06)(251.69,339.94) -- (256.02,346.67)(260.33,334.33) -- (264.75,341)(268.44,328.71) -- (273.15,335.18)(275.89,322.92) -- (281.04,329.04)(282.59,316.59) -- (288.5,321.99)(287.62,309.4) -- (294.89,312.73)(289.34,301.08) -- (297.33,301.56)(290.86,290.72) -- (298.72,292.2)(292.36,281.03) -- (300.28,282.21)(294.21,270.71) -- (302.01,272.47)(297.6,260.38) -- (304.8,263.87)(306.23,251.42) -- (309.27,258.81)(317.69,250.4) -- (317.68,258.4)(327.75,250.42) -- (327.7,258.42)(337.8,250.59) -- (337.54,258.58)(347.68,250.64) -- (347.67,258.64)(357.73,250.76) -- (357.58,258.76)(367.98,251.02) -- (367.72,259.02)(377.42,251.04) -- (377.95,259.02)(385.43,249.08) -- (389.15,256.17)(391.51,243.55) -- (397.96,248.28)(396.13,235.41) -- (403.4,238.75)(399.69,226.51) -- (407.3,228.97)(399.96,219.79) -- (407.37,216.78)(394.75,215.12) -- (398.2,207.9)(386.01,212.4) -- (387.64,204.57)(376.41,210.92) -- (377.32,202.97)(366.64,210) -- (367.26,202.03)(356.86,209.29) -- (357.44,201.31)(346.94,208.49) -- (347.67,200.52)(336.67,207.31) -- (337.81,199.39)(326.51,205.39) -- (328.44,197.62)(316.08,201.74) -- (319.56,194.54)(306.86,194.82) -- (313.07,189.77)(302.95,182.91) -- (310.94,183.33)(306.68,171.74) -- (313.52,175.89)(313.93,163.2) -- (319.15,169.26)(325.59,158.6) -- (325.47,166.6)(336.38,160.13) -- (334.67,167.95)(346.21,162.57) -- (344.15,170.3)(355.89,165.19) -- (353.81,172.91)(365.84,167.79) -- (363.9,175.55)(375.27,169.99) -- (373.59,177.81)(384.94,171.83) -- (383.64,179.73)(394.68,173.12) -- (393.91,181.08)(404.34,173.65) -- (404.23,181.64)(414.01,173.46) -- (414.44,181.45)(423.74,172.59) -- (424.74,180.52)(433.24,171.03) -- (434.82,178.87)(442.54,168.79) -- (444.71,176.49)(451.64,165.86) -- (454.37,173.38)(460.73,162.22) -- (463.87,169.58)(470.34,158.51) -- (473.05,166.04)(479.74,155.16) -- (482.41,162.7)(489.17,151.84) -- (491.81,159.39)(498.6,148.54) -- (501.23,156.1)(508.31,145.18) -- (510.92,152.74)(517.48,142.02) -- (520.08,149.59)(527.25,138.68) -- (529.82,146.26)(536.47,135.56) -- (539.02,143.14) ;
%Curve Lines [id:da3876903888663663] 
\draw [color={rgb, 255:red, 155; green, 155; blue, 155 }  ,draw opacity=1 ]   (201.2,205.6) .. controls (202.36,240.13) and (263.37,280.12) .. (284.37,284.31) ;
\draw [shift={(286.2,284.6)}, rotate = 186.01] [color={rgb, 255:red, 155; green, 155; blue, 155 }  ,draw opacity=1 ][line width=0.75]    (10.93,-3.29) .. controls (6.95,-1.4) and (3.31,-0.3) .. (0,0) .. controls (3.31,0.3) and (6.95,1.4) .. (10.93,3.29)   ;
%Curve Lines [id:da04910207175146397] 
\draw [color={rgb, 255:red, 155; green, 155; blue, 155 }  ,draw opacity=1 ]   (417.2,118.6) .. controls (391.98,144.79) and (362.06,185.09) .. (367.58,203) ;
\draw [shift={(368.2,204.6)}, rotate = 244.8] [color={rgb, 255:red, 155; green, 155; blue, 155 }  ,draw opacity=1 ][line width=0.75]    (10.93,-3.29) .. controls (6.95,-1.4) and (3.31,-0.3) .. (0,0) .. controls (3.31,0.3) and (6.95,1.4) .. (10.93,3.29)   ;
%Shape: Circle [id:dp48139133948616] 
\draw   (311.5,254.4) .. controls (311.5,240.59) and (322.69,229.4) .. (336.5,229.4) .. controls (350.31,229.4) and (361.5,240.59) .. (361.5,254.4) .. controls (361.5,268.21) and (350.31,279.4) .. (336.5,279.4) .. controls (322.69,279.4) and (311.5,268.21) .. (311.5,254.4) -- cycle ;
%Curve Lines [id:da6583539286100435] 
\draw    (348.2,264.6) .. controls (357.08,269.04) and (449.39,344.06) .. (466.63,356.5) ;
\draw [shift={(468.2,357.6)}, rotate = 213.69] [color={rgb, 255:red, 0; green, 0; blue, 0 }  ][line width=0.75]    (10.93,-3.29) .. controls (6.95,-1.4) and (3.31,-0.3) .. (0,0) .. controls (3.31,0.3) and (6.95,1.4) .. (10.93,3.29)   ;

% Text Node
\draw (559,272.4) node [anchor=north west][inner sep=0.75pt]    {$\mathbb{R}^{a}$};
% Text Node
\draw (298,77.4) node [anchor=north west][inner sep=0.75pt]    {$\mathbb{R}^{b}$};
% Text Node
\draw (415,87.4) node [anchor=north west][inner sep=0.75pt]    {$h\left(\mathbb{R}^{a}\right)$};
% Text Node
\draw (165,174.4) node [anchor=north west][inner sep=0.75pt]    {$\overline{h}_{1}\left(\mathbb{R}^{a}\right)$};
% Text Node
\draw (224,306.4) node [anchor=north west][inner sep=0.75pt]    {$h( E)$};
% Text Node
\draw (476,350.4) node [anchor=north west][inner sep=0.75pt]    {$E_{1} =h( E_{1})$};


\end{tikzpicture}
}}

\newcommand{\FigSixOne}{%
  \resizebox{.84\textwidth}{!}{\tikzset{every picture/.style={line width=0.75pt}} %set default line width to 0.75pt

\begin{tikzpicture}[x=0.75pt,y=0.75pt,yscale=-1,xscale=1]
%uncomment if require: \path (0,724); %set diagram left start at 0, and has height of 724

%Curve Lines [id:da9620044273585225]
\draw    (140.2,413.8) .. controls (326.2,-6.2) and (486.2,699.8) .. (600.2,304.8) ;
%Curve Lines [id:da10895257227162702]
\draw    (52.2,386.8) .. controls (93.2,294.8) and (126.2,259.8) .. (140.2,263.8) ;
%Straight Lines [id:da8327014062778634]
\draw    (140.2,263.8) -- (273.2,274.8) ;
%Curve Lines [id:da4986053383299296]
\draw [color={rgb, 255:red, 128; green, 128; blue, 128 }  ,draw opacity=1 ]   (120.2,361.8) .. controls (161.2,269.8) and (192.7,267.9) .. (206.7,269.3) ;
%Straight Lines [id:da9571954478386947]
\draw    (52.2,386.8) -- (140.2,413.8) ;
%Curve Lines [id:da8062034220654326]
\draw    (30.2,366.8) .. controls (51.2,364.8) and (93.2,361.8) .. (120.2,361.8) ;
%Curve Lines [id:da40748546279172926]
\draw [color={rgb, 255:red, 128; green, 128; blue, 128 }  ,draw opacity=1 ] [dash pattern={on 4.5pt off 4.5pt}]  (120.2,361.8) .. controls (136.4,361.2) and (148.8,361.6) .. (164.2,361.8) ;
%Curve Lines [id:da8395773850437817]
\draw  [dash pattern={on 4.5pt off 4.5pt}]  (140.2,263.8) .. controls (166.2,274.8) and (196.2,291.8) .. (210.2,301.8) ;
%Curve Lines [id:da9403563547803104]
\draw    (210.2,301.8) .. controls (232.2,316.8) and (297.2,386.8) .. (315.2,403.8) .. controls (333.2,420.8) and (347.2,435.8) .. (365.2,449.8) ;
%Curve Lines [id:da21283109640207787]
\draw    (365.2,449.8) .. controls (388.2,461.8) and (492.4,454.9) .. (532.4,424.9) ;
%Curve Lines [id:da905844513019243]
\draw  [dash pattern={on 4.5pt off 4.5pt}]  (365.2,449.8) .. controls (376.2,446.8) and (409.2,441.8) .. (456.2,403.8) ;
%Curve Lines [id:da24023412254134335]
\draw    (456.2,403.8) .. controls (496.2,373.8) and (497.2,346.8) .. (514.2,297.8) ;
%Curve Lines [id:da617810943510469]
\draw    (514.2,297.8) .. controls (538.2,302.8) and (574.2,306.8) .. (600.2,304.8) ;
%Curve Lines [id:da05339788381467914]
\draw [color={rgb, 255:red, 128; green, 128; blue, 128 }  ,draw opacity=1 ]   (164.2,361.8) .. controls (205.2,360.8) and (310.2,363.8) .. (349.2,367.8) ;
%Curve Lines [id:da37032907704300666]
\draw  [dash pattern={on 4.5pt off 4.5pt}]  (349.2,367.8) .. controls (364.2,369.8) and (405.92,375.6) .. (427.2,378.8) ;
%Curve Lines [id:da5858810449257141]
\draw    (427.2,378.8) .. controls (458.2,384.8) and (484.2,390.8) .. (513.2,396.8) ;
%Curve Lines [id:da5792510662820718]
\draw  [dash pattern={on 4.5pt off 4.5pt}]  (513.2,396.8) .. controls (527.12,400) and (539.6,402.9) .. (553.52,406.8) ;
%Curve Lines [id:da04832578340355309]
\draw    (553.52,406.8) .. controls (575.3,413.3) and (588.3,416.8) .. (612.2,425.2) ;
%Curve Lines [id:da6743216618309311]
\draw [color={rgb, 255:red, 128; green, 128; blue, 128 }  ,draw opacity=1 ]   (206.7,269.3) .. controls (217.5,271.4) and (226.5,279.4) .. (232.7,286.4) ;
%Curve Lines [id:da6352281891905399]
\draw [color={rgb, 255:red, 128; green, 128; blue, 128 }  ,draw opacity=1 ]   (232.7,286.4) .. controls (279.3,327.4) and (330.8,358.9) .. (349.2,367.8) ;
%Shape: Circle [id:dp256915589824944]
\draw  [fill={rgb, 255:red, 0; green, 0; blue, 0 }  ,fill opacity=1 ] (118.1,361.8) .. controls (118.1,360.64) and (119.04,359.7) .. (120.2,359.7) .. controls (121.36,359.7) and (122.3,360.64) .. (122.3,361.8) .. controls (122.3,362.96) and (121.36,363.9) .. (120.2,363.9) .. controls (119.04,363.9) and (118.1,362.96) .. (118.1,361.8) -- cycle ;
%Shape: Circle [id:dp662684538575718]
\draw  [fill={rgb, 255:red, 0; green, 0; blue, 0 }  ,fill opacity=1 ] (347.1,367.8) .. controls (347.1,366.64) and (348.04,365.7) .. (349.2,365.7) .. controls (350.36,365.7) and (351.3,366.64) .. (351.3,367.8) .. controls (351.3,368.96) and (350.36,369.9) .. (349.2,369.9) .. controls (348.04,369.9) and (347.1,368.96) .. (347.1,367.8) -- cycle ;

% Text Node
\draw (114.9,366.4) node [anchor=north west][inner sep=0.75pt]    {$p$};
% Text Node
\draw (351.2,373.3) node [anchor=north west][inner sep=0.75pt]    {$q$};
% Text Node
\draw (287.4,344.4) node [anchor=north west][inner sep=0.75pt]  [color={rgb, 255:red, 128; green, 128; blue, 128 }  ,opacity=1 ]  {$L$};
% Text Node
\draw (335,453.4) node [anchor=north west][inner sep=0.75pt]    {$M$};
% Text Node
\draw (614.2,428.6) node [anchor=north west][inner sep=0.75pt]    {$M'$};

% 收紧图形边界框；不裁切也不改动任何曲线，只去除控制点造成的上下空白。
\pgfresetboundingbox
\path[use as bounding box] (20,245) rectangle (670,475);

\end{tikzpicture}
}}

\newcommand{\FigSixTwo}{%
  \resizebox{.78\textwidth}{!}{% Pattern Info
 
\tikzset{
pattern size/.store in=\mcSize, 
pattern size = 5pt,
pattern thickness/.store in=\mcThickness, 
pattern thickness = 0.3pt,
pattern radius/.store in=\mcRadius, 
pattern radius = 1pt}
\makeatletter
\pgfutil@ifundefined{pgf@pattern@name@_26kpy2v0r}{
\pgfdeclarepatternformonly[\mcThickness,\mcSize]{_26kpy2v0r}
{\pgfqpoint{0pt}{0pt}}
{\pgfpoint{\mcSize}{\mcSize}}
{\pgfpoint{\mcSize}{\mcSize}}
{
\pgfsetcolor{\tikz@pattern@color}
\pgfsetlinewidth{\mcThickness}
\pgfpathmoveto{\pgfqpoint{0pt}{\mcSize}}
\pgfpathlineto{\pgfpoint{\mcSize+\mcThickness}{-\mcThickness}}
\pgfpathmoveto{\pgfqpoint{0pt}{0pt}}
\pgfpathlineto{\pgfpoint{\mcSize+\mcThickness}{\mcSize+\mcThickness}}
\pgfusepath{stroke}
}}
\makeatother
\tikzset{every picture/.style={line width=0.75pt}} %set default line width to 0.75pt        

\begin{tikzpicture}[x=0.75pt,y=0.75pt,yscale=-1,xscale=1]
%uncomment if require: \path (0,724); %set diagram left start at 0, and has height of 724

%Shape: Ellipse [id:dp12751974195854354] 
\draw  [pattern=_26kpy2v0r,pattern size=12pt,pattern thickness=0.22pt,pattern radius=0pt, pattern color={rgb, 255:red, 215; green, 215; blue, 215}] (100,371.8) .. controls (100,290.72) and (196.75,225) .. (316.1,225) .. controls (435.45,225) and (532.2,290.72) .. (532.2,371.8) .. controls (532.2,452.88) and (435.45,518.6) .. (316.1,518.6) .. controls (196.75,518.6) and (100,452.88) .. (100,371.8) -- cycle ;
%Shape: Parabola [id:dp7097132516162088] 
\draw   (584.7,459.6) .. controls (406.53,205.73) and (228.37,205.73) .. (50.2,459.6) ;
%Shape: Parabola [id:dp18585049318765634] 
\draw   (48.85,270) .. controls (227.02,523.87) and (405.18,523.87) .. (583.35,270) ;

% Text Node
%Straight Lines [id:da6462481686425137]
\draw    (317.45,269.2) ;
\draw [shift={(317.45,269.2)}, rotate = 180] [fill={rgb, 255:red, 0; green, 0; blue, 0 }  ][line width=0.08]  [draw opacity=0] (8.93,-4.29) -- (0,0) -- (8.93,4.29) -- cycle    ;
%Straight Lines [id:da8521475132911274]
\draw    (316.1,460.4) ;
\draw [shift={(316.1,460.4)}, rotate = 180] [fill={rgb, 255:red, 0; green, 0; blue, 0 }  ][line width=0.08]  [draw opacity=0] (8.93,-4.29) -- (0,0) -- (8.93,4.29) -- cycle    ;

\draw (146,357.4) node [anchor=north west][inner sep=0.75pt]    {$a$};
% Text Node
\draw (481,357.4) node [anchor=north west][inner sep=0.75pt]    {$b$};
% Text Node
\draw (311,466.4) node [anchor=north west][inner sep=0.75pt]    {$C_{0}$};
% Text Node
\draw (300,244.4) node [anchor=north west][inner sep=0.75pt]    {$C'_{0}$};
% Text Node
\draw (194,266.4) node [anchor=north west][inner sep=0.75pt]    {$U$};

\end{tikzpicture}
}}

\newcommand{\FigSixThree}{%
  \resizebox{.78\textwidth}{!}{% Pattern Info

\tikzset{
pattern size/.store in=\mcSize,
pattern size = 5pt,
pattern thickness/.store in=\mcThickness,
pattern thickness = 0.3pt,
pattern radius/.store in=\mcRadius,
pattern radius = 1pt}
\makeatletter
\pgfutil@ifundefined{pgf@pattern@name@_bni9gw1ql}{
\pgfdeclarepatternformonly[\mcThickness,\mcSize]{_bni9gw1ql}
{\pgfqpoint{0pt}{0pt}}
{\pgfpoint{\mcSize}{\mcSize}}
{\pgfpoint{\mcSize}{\mcSize}}
{
\pgfsetcolor{\tikz@pattern@color}
\pgfsetlinewidth{\mcThickness}
\pgfpathmoveto{\pgfqpoint{0pt}{\mcSize}}
\pgfpathlineto{\pgfpoint{\mcSize+\mcThickness}{-\mcThickness}}
\pgfpathmoveto{\pgfqpoint{0pt}{0pt}}
\pgfpathlineto{\pgfpoint{\mcSize+\mcThickness}{\mcSize+\mcThickness}}
\pgfusepath{stroke}
}}
\makeatother
\tikzset{every picture/.style={line width=0.75pt}} %set default line width to 0.75pt

\begin{tikzpicture}[x=0.75pt,y=0.75pt,yscale=-1,xscale=1]
%uncomment if require: \path (0,724); %set diagram left start at 0, and has height of 724

%Shape: Ellipse [id:dp12751974195854354]
\draw  [pattern=_bni9gw1ql,pattern size=12pt,pattern thickness=0.22pt,pattern radius=0pt, pattern color={rgb, 255:red, 215; green, 215; blue, 215}] (100,371.8) .. controls (100,290.72) and (196.75,225) .. (316.1,225) .. controls (435.45,225) and (532.2,290.72) .. (532.2,371.8) .. controls (532.2,452.88) and (435.45,518.6) .. (316.1,518.6) .. controls (196.75,518.6) and (100,452.88) .. (100,371.8) -- cycle ;
%Shape: Parabola [id:dp7097132516162088]
\draw   (584.7,459.6) .. controls (406.53,205.73) and (228.37,205.73) .. (50.2,459.6) ;
%Shape: Parabola [id:dp18585049318765634]
\draw   (48.85,270) .. controls (227.02,523.87) and (405.18,523.87) .. (583.35,270) ;

% 以下五条样条以 x=316.1 为严格对称轴；左半边保留用户坐标，
% 右半边的端点和 Bézier 控制点均取其镜像。
%Curve Lines [id:da7533471946150027]
\draw [color={rgb, 255:red, 139; green, 87; blue, 42 },draw opacity=1]
  (237.07,443.33)
  .. controls (249.23,452.47) and (284.1,432.67) .. (316.1,433.33)
  .. controls (348.1,432.67) and (382.97,452.47) .. (395.13,443.33);
%Curve Lines [id:da6025026154144115]
\draw [color={rgb, 255:red, 139; green, 87; blue, 42 },draw opacity=1]
  (185.73,414.67)
  .. controls (220.4,440) and (279.57,392.33) .. (316.1,391.8)
  .. controls (352.63,392.33) and (411.8,440) .. (446.47,414.67);
%Curve Lines [id:da7505325398863362]
\draw [color={rgb, 255:red, 139; green, 87; blue, 42 },draw opacity=1]
  (158.4,394.67)
  .. controls (193.07,420) and (279.57,352.33) .. (316.1,351.8)
  .. controls (352.63,352.33) and (439.13,420) .. (473.8,394.67);
%Curve Lines [id:da7650080957487818]
\draw [color={rgb, 255:red, 139; green, 87; blue, 42 },draw opacity=1]
  (135.07,373.33)
  .. controls (152.4,386) and (184.04,351.77) .. (220.4,330.67)
  .. controls (256.76,309.57) and (297.83,301.6) .. (316.1,301.33)
  .. controls (334.37,301.6) and (375.44,309.57) .. (411.8,330.67)
  .. controls (448.16,351.77) and (479.8,386) .. (497.13,373.33);
%Curve Lines [id:da18903048061491867]
\draw [color={rgb, 255:red, 139; green, 87; blue, 42 },draw opacity=1]
  (113.07,350)
  .. controls (121.73,356.33) and (130.94,341.73) .. (151.07,325.33)
  .. controls (171.19,308.94) and (199.55,291.22) .. (217.73,280.67)
  .. controls (235.91,270.12) and (296.13,253.6) .. (316.1,253.33)
  .. controls (336.07,253.6) and (396.29,270.12) .. (414.47,280.67)
  .. controls (432.65,291.22) and (461.01,308.94) .. (481.13,325.33)
  .. controls (501.26,341.73) and (510.47,356.33) .. (519.13,350);

% Text Node
\draw (318.1,463.8) node [anchor=north west][inner sep=0.75pt]    {$C_{0}$};
% Text Node
\draw (305.43,275.13) node [anchor=north west][inner sep=0.75pt]    {$C'_{0}$};
% Text Node
\draw (271.2,233.93) node [anchor=north west][inner sep=0.75pt]  [color={rgb, 255:red, 139; green, 87; blue, 42 }  ,opacity=1 ]  {$G_{1}( U\cap C_{0})$};

\end{tikzpicture}
}}

\newcommand{\FigSixFour}{%
  \resizebox{.90\textwidth}{!}{\tikzset{every picture/.style={line width=0.75pt}}

\begin{tikzpicture}[x=0.75pt,y=0.75pt,yscale=-1,xscale=1]
%Curve Lines [id:da5513676269251793]
\draw (40.2,293.4) .. controls (80.2,263.4) and (97.2,242.4) .. (137.2,212.4) .. controls (177.2,182.4) and (210.2,173.4) .. (234.2,170.4) .. controls (258.2,167.4) and (306.2,169.4) .. (327.2,174.4) .. controls (348.2,179.4) and (361.2,187.4) .. (397.2,202.4) .. controls (433.2,217.4) and (475.2,243.4) .. (501.2,255.4) .. controls (527.2,267.4) and (600.2,293.4) .. (611.2,315.4) ;
%Curve Lines [id:da7211046777786992]
\draw [color={rgb, 255:red, 74; green, 144; blue, 226},draw opacity=1] (56.2,307.4) .. controls (95.2,270.4) and (108.2,269.4) .. (140.2,235.4) .. controls (172.2,201.4) and (212.2,190.4) .. (236.2,187.4) .. controls (260.2,184.4) and (308.2,190.4) .. (318.2,192.4) .. controls (328.2,194.4) and (381.2,232.4) .. (397.2,239.4) .. controls (413.2,246.4) and (470.2,251.4) .. (496.2,263.4) .. controls (522.2,275.4) and (594.2,303.4) .. (606.2,321.4) ;
%Curve Lines [id:da3249325297243797]
\draw (72.2,323.4) .. controls (112.2,293.4) and (127.2,265.4) .. (167.2,235.4) .. controls (207.2,205.4) and (233.2,207.4) .. (255.2,207.4) .. controls (277.2,207.4) and (302.2,211.4) .. (321.2,217.4) .. controls (340.2,223.4) and (372.2,243.4) .. (393.2,251.4) .. controls (414.2,259.4) and (445.2,274.4) .. (468.2,278.4) .. controls (491.2,282.4) and (503.2,288.77) .. (529.2,298.4) .. controls (555.2,308.02) and (586.2,327.9) .. (596.2,338.4) ;
%Curve Lines [id:da7933953701836717]
\draw (34.2,128.4) .. controls (79.2,119.4) and (114.2,153.4) .. (124.2,172.4) .. controls (134.2,191.4) and (179.2,283.4) .. (194.2,295.4) .. controls (209.2,307.4) and (237.2,322.4) .. (253.2,326.4) .. controls (269.2,330.4) and (308.2,336.4) .. (335.2,340.4) .. controls (362.2,344.4) and (389.2,348.4) .. (420.2,342.4) .. controls (451.2,336.4) and (484.82,328.52) .. (506.2,313.4) .. controls (527.57,298.27) and (559.7,246.9) .. (572.2,234.4) .. controls (584.7,221.9) and (612.2,183.4) .. (656.2,194.4) ;
%Curve Lines [id:da05584228564866156]
\draw [color={rgb, 255:red, 208; green, 2; blue, 27},draw opacity=1] (32.7,147.4) .. controls (70.2,127.9) and (105.6,175.4) .. (115.6,194.4) .. controls (125.6,213.4) and (170.6,305.4) .. (185.6,317.4) .. controls (200.6,329.4) and (228.6,344.4) .. (244.6,348.4) .. controls (260.6,352.4) and (299.6,358.4) .. (326.6,362.4) .. controls (353.6,366.4) and (380.6,370.4) .. (411.6,364.4) .. controls (442.6,358.4) and (476.23,350.52) .. (497.6,335.4) .. controls (518.98,320.27) and (549.7,280.9) .. (562.2,268.4) .. controls (574.7,255.9) and (614.8,205.4) .. (653.2,206.4) ;
%Curve Lines [id:da8356231429288252]
\draw (30.6,168.4) .. controls (70.6,138.4) and (100.6,190.4) .. (110.6,209.4) .. controls (120.6,228.4) and (165.6,320.4) .. (180.6,332.4) .. controls (195.6,344.4) and (224.6,358.4) .. (237.2,364.4) .. controls (249.8,370.4) and (290.2,379.4) .. (317.2,383.4) .. controls (344.2,387.4) and (399.2,386.4) .. (420.2,382.4) .. controls (441.2,378.4) and (488.2,364.4) .. (502.2,356.4) .. controls (516.2,348.4) and (545.2,318.4) .. (553.2,310.4) .. controls (561.2,302.4) and (617.2,244.4) .. (650.2,229.4) ;
%Shape: Ellipse [id:dp786521588870946]
\draw (105.2,235.4) .. controls (105.2,211.93) and (120.87,192.9) .. (140.2,192.9) .. controls (159.53,192.9) and (175.2,211.93) .. (175.2,235.4) .. controls (175.2,258.87) and (159.53,277.9) .. (140.2,277.9) .. controls (120.87,277.9) and (105.2,258.87) .. (105.2,235.4) -- cycle ;
%Shape: Ellipse [id:dp8023382049051939]
\draw (75.8,235.4) .. controls (75.8,198.29) and (104.63,168.2) .. (140.2,168.2) .. controls (175.77,168.2) and (204.6,198.29) .. (204.6,235.4) .. controls (204.6,272.51) and (175.77,302.6) .. (140.2,302.6) .. controls (104.63,302.6) and (75.8,272.51) .. (75.8,235.4) -- cycle ;
%Shape: Ellipse [id:dp6345230646685797]
\draw (177.1,194.4) .. controls (177.1,169.55) and (205.1,149.4) .. (239.65,149.4) .. controls (274.2,149.4) and (302.2,169.55) .. (302.2,194.4) .. controls (302.2,219.25) and (274.2,239.4) .. (239.65,239.4) .. controls (205.1,239.4) and (177.1,219.25) .. (177.1,194.4) -- cycle ;
%Shape: Ellipse [id:dp7098931107611396]
\draw (270.6,195.15) .. controls (270.6,161.05) and (298.94,133.4) .. (333.9,133.4) .. controls (368.86,133.4) and (397.2,161.05) .. (397.2,195.15) .. controls (397.2,229.25) and (368.86,256.9) .. (333.9,256.9) .. controls (298.94,256.9) and (270.6,229.25) .. (270.6,195.15) -- cycle ;
%Shape: Ellipse [id:dp12966068998368774]
\draw (360.93,211.74) .. controls (369.75,187.11) and (405.71,177.45) .. (441.26,190.17) .. controls (476.82,202.9) and (498.49,233.18) .. (489.68,257.81) .. controls (480.86,282.44) and (444.89,292.1) .. (409.34,279.38) .. controls (373.79,266.66) and (352.12,236.37) .. (360.93,211.74) -- cycle ;
%Shape: Ellipse [id:dp7775856234119033]
\draw (146.83,298.83) .. controls (155.03,275.9) and (190.5,267.63) .. (226.06,280.35) .. controls (261.61,293.07) and (283.78,321.97) .. (275.57,344.9) .. controls (267.37,367.83) and (231.89,376.1) .. (196.34,363.38) .. controls (160.79,350.66) and (138.62,321.75) .. (146.83,298.83) -- cycle ;
%Shape: Ellipse [id:dp864059912702568]
\draw (241,350.9) .. controls (241,326.32) and (265.22,306.4) .. (295.1,306.4) .. controls (324.98,306.4) and (349.2,326.32) .. (349.2,350.9) .. controls (349.2,375.48) and (324.98,395.4) .. (295.1,395.4) .. controls (265.22,395.4) and (241,375.48) .. (241,350.9) -- cycle ;
%Shape: Ellipse [id:dp8528607753051314]
\draw (305,358.2) .. controls (305,331.58) and (343.55,310) .. (391.1,310) .. controls (438.65,310) and (477.2,331.58) .. (477.2,358.2) .. controls (477.2,384.82) and (438.65,406.4) .. (391.1,406.4) .. controls (343.55,406.4) and (305,384.82) .. (305,358.2) -- cycle ;
%Shape: Ellipse [id:dp4768239619464324]
\draw [fill={rgb, 255:red, 155; green, 155; blue, 155},fill opacity=0.11] (446.21,293.6) .. controls (452.54,245.49) and (490.62,210.83) .. (531.27,216.18) .. controls (571.91,221.53) and (599.72,264.87) .. (593.39,312.97) .. controls (587.05,361.08) and (548.97,395.74) .. (508.33,390.39) .. controls (467.69,385.04) and (439.88,341.7) .. (446.21,293.6) -- cycle ;
%Shape: Ellipse [id:dp8685388596017811]
\draw [fill={rgb, 255:red, 155; green, 155; blue, 155},fill opacity=0.4] (504.4,298.02) .. controls (501.96,278.58) and (518.5,260.5) .. (541.33,257.63) .. controls (564.17,254.77) and (584.66,268.19) .. (587.1,287.63) .. controls (589.54,307.06) and (573.01,325.14) .. (550.17,328.01) .. controls (527.34,330.88) and (506.84,317.45) .. (504.4,298.02) -- cycle ;
%Shape: Circle [id:dp856308539569508]
\draw [fill={rgb, 255:red, 0; green, 0; blue, 0},fill opacity=1] (545.6,286.51) .. controls (545.6,285.75) and (546.22,285.13) .. (546.98,285.13) .. controls (547.74,285.13) and (548.36,285.75) .. (548.36,286.51) .. controls (548.36,287.27) and (547.74,287.89) .. (546.98,287.89) .. controls (546.22,287.89) and (545.6,287.27) .. (545.6,286.51) -- cycle ;
%Straight Lines [id:da967030685461137]
\draw (25,149.9) -- (26.4,166.4) ;
\draw [shift={(26.4,166.4)},rotate=265.15,color={rgb, 255:red, 0; green, 0; blue, 0},line width=0.75] (0,5.59) -- (0,-5.59) ;
\draw [shift={(25,149.9)},rotate=265.15,color={rgb, 255:red, 0; green, 0; blue, 0},line width=0.75] (0,5.59) -- (0,-5.59) ;
%Straight Lines [id:da9025702377469016]
\draw (38,297.4) -- (50,311.4) ;
\draw [shift={(50,311.4)},rotate=229.4,color={rgb, 255:red, 0; green, 0; blue, 0},line width=0.75] (0,5.59) -- (0,-5.59) ;
\draw [shift={(38,297.4)},rotate=229.4,color={rgb, 255:red, 0; green, 0; blue, 0},line width=0.75] (0,5.59) -- (0,-5.59) ;

% Text Nodes
\draw (7.5,124.97) node [anchor=north west,inner sep=0.75pt] {$\textcolor[rgb]{0.82,0.01,0.11}{M}$};
\draw (46.5,315.47) node [anchor=north west,inner sep=0.75pt] {$\textcolor[rgb]{0.29,0.56,0.89}{M'}$};
\draw (508.2,365.97) node [anchor=north west,inner sep=0.75pt] {$W_{i}$};
\draw (508.2,274.97) node [anchor=north west,inner sep=0.75pt] {$N_{i}$};
\draw (545.32,292.73) node [anchor=north west,inner sep=0.75pt] {$p_{i}$};
\draw (5.5,150.9) node [anchor=north west,inner sep=0.75pt] {$T$};
\draw (20.5,304.4) node [anchor=north west,inner sep=0.75pt] {$T'$};

\end{tikzpicture}
}}

\newcommand{\FigSixFive}{%
  \resizebox{.78\textwidth}{!}{% Pattern Info

\tikzset{
pattern size/.store in=\mcSize,
pattern size = 5pt,
pattern thickness/.store in=\mcThickness,
pattern thickness = 0.3pt,
pattern radius/.store in=\mcRadius,
pattern radius = 1pt}
\makeatletter
\pgfutil@ifundefined{pgf@pattern@name@_bgrd8j6uq}{
\pgfdeclarepatternformonly[\mcThickness,\mcSize]{_bgrd8j6uq}
{\pgfqpoint{0pt}{0pt}}
{\pgfpoint{\mcSize}{\mcSize}}
{\pgfpoint{\mcSize}{\mcSize}}
{
\pgfsetcolor{\tikz@pattern@color}
\pgfsetlinewidth{\mcThickness}
\pgfpathmoveto{\pgfqpoint{0pt}{\mcSize}}
\pgfpathlineto{\pgfpoint{\mcSize+\mcThickness}{-\mcThickness}}
\pgfpathmoveto{\pgfqpoint{0pt}{0pt}}
\pgfpathlineto{\pgfpoint{\mcSize+\mcThickness}{\mcSize+\mcThickness}}
\pgfusepath{stroke}
}}
\makeatother
\tikzset{every picture/.style={line width=0.75pt}} %set default line width to 0.75pt

\begin{tikzpicture}[x=0.75pt,y=0.75pt,yscale=-1,xscale=1]
%uncomment if require: \path (0,724); %set diagram left start at 0, and has height of 724

%Shape: Ellipse [id:dp12751974195854354]
\draw  [pattern=_bgrd8j6uq,pattern size=12pt,pattern thickness=0.22pt,pattern radius=0pt, pattern color={rgb, 255:red, 215; green, 215; blue, 215}] (100,371.8) .. controls (100,290.72) and (196.75,225) .. (316.1,225) .. controls (435.45,225) and (532.2,290.72) .. (532.2,371.8) .. controls (532.2,452.88) and (435.45,518.6) .. (316.1,518.6) .. controls (196.75,518.6) and (100,452.88) .. (100,371.8) -- cycle ;
%Shape: Parabola [id:dp7097132516162088]
\draw   (584.7,459.6) .. controls (406.53,205.73) and (228.37,205.73) .. (50.2,459.6) ;
%Shape: Parabola [id:dp18585049318765634]
\draw   (48.85,270) .. controls (227.02,523.87) and (405.18,523.87) .. (583.35,270) ;
%Shape: Ellipse [id:dp6935180515795655]
\draw  [fill={rgb, 255:red, 255; green, 255; blue, 255 }  ,fill opacity=1 ] (236.5,371.8) .. controls (236.5,346.78) and (272.14,326.5) .. (316.1,326.5) .. controls (360.06,326.5) and (395.7,346.78) .. (395.7,371.8) .. controls (395.7,396.82) and (360.06,417.1) .. (316.1,417.1) .. controls (272.14,417.1) and (236.5,396.82) .. (236.5,371.8) -- cycle ;


% 原图 6.5 中沿开口朝上的抛物线设置的法向箭头。
\draw[->] (156,392)--(169,375);
\draw[->] (193,420)--(205,399);
\draw[->] (231,442)--(240,418);
\draw[->] (268,455)--(273,429);
\draw[->] (300,460)--(301,432);
\draw[->] (332,460)--(331,432);
\draw[->] (364,454)--(359,428);
\draw[->] (401,442)--(391,418);
\draw[->] (438,420)--(426,400);
\draw[->] (477,392)--(464,375);

% Text Node
%Straight Lines [id:da6462481686425137]
\draw    (317.45,269.2) ;
\draw [shift={(317.45,269.2)}, rotate = 180] [fill={rgb, 255:red, 0; green, 0; blue, 0 }  ][line width=0.08]  [draw opacity=0] (8.93,-4.29) -- (0,0) -- (8.93,4.29) -- cycle    ;
%Straight Lines [id:da8521475132911274]
\draw    (316.1,460.4) ;
\draw [shift={(316.1,460.4)}, rotate = 180] [fill={rgb, 255:red, 0; green, 0; blue, 0 }  ][line width=0.08]  [draw opacity=0] (8.93,-4.29) -- (0,0) -- (8.93,4.29) -- cycle    ;

\draw (146,357.4) node [anchor=north west][inner sep=0.75pt]    {$a$};
% Text Node
\draw (481,357.4) node [anchor=north west][inner sep=0.75pt]    {$b$};
% Text Node
\draw (311,466.4) node [anchor=north west][inner sep=0.75pt]    {$C_{0}$};
% Text Node
\draw (300,244.4) node [anchor=north west][inner sep=0.75pt]    {$C'_{0}$};
% Text Node
\draw (194,266.4) node [anchor=north west][inner sep=0.75pt]    {$N$};
% Text Node
\draw (325,337.4) node [anchor=north west][inner sep=0.75pt]    {$S$};

\end{tikzpicture}
}}

\newcommand{\FigSixSix}{%
  \resizebox{.78\textwidth}{!}{% Pattern Info

\tikzset{
pattern size/.store in=\mcSize,
pattern size = 5pt,
pattern thickness/.store in=\mcThickness,
pattern thickness = 0.3pt,
pattern radius/.store in=\mcRadius,
pattern radius = 1pt}
\makeatletter
\pgfutil@ifundefined{pgf@pattern@name@_xdavnpld8}{
\pgfdeclarepatternformonly[\mcThickness,\mcSize]{_xdavnpld8}
{\pgfqpoint{0pt}{0pt}}
{\pgfpoint{\mcSize}{\mcSize}}
{\pgfpoint{\mcSize}{\mcSize}}
{
\pgfsetcolor{\tikz@pattern@color}
\pgfsetlinewidth{\mcThickness}
\pgfpathmoveto{\pgfqpoint{0pt}{\mcSize}}
\pgfpathlineto{\pgfpoint{\mcSize+\mcThickness}{-\mcThickness}}
\pgfpathmoveto{\pgfqpoint{0pt}{0pt}}
\pgfpathlineto{\pgfpoint{\mcSize+\mcThickness}{\mcSize+\mcThickness}}
\pgfusepath{stroke}
}}
\makeatother
\tikzset{every picture/.style={line width=0.75pt}} %set default line width to 0.75pt

\begin{tikzpicture}[x=0.75pt,y=0.75pt,yscale=-1,xscale=1]
%uncomment if require: \path (0,724); %set diagram left start at 0, and has height of 724

%Shape: Ellipse [id:dp12751974195854354]
\draw  [pattern=_xdavnpld8,pattern size=12pt,pattern thickness=0.22pt,pattern radius=0pt, pattern color={rgb, 255:red, 215; green, 215; blue, 215}] (100,371.8) .. controls (100,290.72) and (196.75,225) .. (316.1,225) .. controls (435.45,225) and (532.2,290.72) .. (532.2,371.8) .. controls (532.2,452.88) and (435.45,518.6) .. (316.1,518.6) .. controls (196.75,518.6) and (100,452.88) .. (100,371.8) -- cycle ;
%Shape: Parabola [id:dp7097132516162088]
\draw   (584.7,459.6) .. controls (406.53,205.73) and (228.37,205.73) .. (50.2,459.6) ;
%Shape: Parabola [id:dp18585049318765634]
\draw   (48.85,270) .. controls (227.02,523.87) and (405.18,523.87) .. (583.35,270) ;
%Straight Lines [id:da8169258507934456]
\draw [line width=1.5]    (127.5,365.5) -- (166.87,404.87) ;
\draw [shift={(169.7,407.7)}, rotate = 225] [fill={rgb, 255:red, 0; green, 0; blue, 0 }  ][line width=0.08]  [draw opacity=0] (11.61,-5.58) -- (0,0) -- (11.61,5.58) -- cycle    ;
%Shape: Right Angle [id:dp7344047988263653]
\draw   (133,371) -- (128.15,375.85) -- (122.65,370.35) ;
%Straight Lines [id:da15089389075617865]
\draw [line width=1.5]    (507,364.4) -- (473.53,397.87) ;
\draw [shift={(470.7,400.7)}, rotate = 315] [fill={rgb, 255:red, 0; green, 0; blue, 0 }  ][line width=0.08]  [draw opacity=0] (11.61,-5.58) -- (0,0) -- (11.61,5.58) -- cycle    ;
%Straight Lines [id:da7410107028706643]
\draw [line width=1.5]    (540.47,330.93) -- (507,364.4) ;
\draw [shift={(543.3,328.1)}, rotate = 135] [fill={rgb, 255:red, 0; green, 0; blue, 0 }  ][line width=0.08]  [draw opacity=0] (11.61,-5.58) -- (0,0) -- (11.61,5.58) -- cycle    ;
%Shape: Right Angle [id:dp22345946357963165]
\draw   (511.75,369.15) -- (507.5,373.4) -- (502.75,368.65) ;
%Shape: Right Angle [id:dp04596191707344943]
\draw   (510.75,360.45) -- (515.7,365.4) -- (511.85,369.25) ;
%Straight Lines [id:da6462481686425137]
\draw    (317.45,269.2) ;
\draw [shift={(317.45,269.2)}, rotate = 180] [fill={rgb, 255:red, 0; green, 0; blue, 0 }  ][line width=0.08]  [draw opacity=0] (8.93,-4.29) -- (0,0) -- (8.93,4.29) -- cycle    ;
%Straight Lines [id:da8521475132911274]
\draw    (316.1,460.4) ;
\draw [shift={(316.1,460.4)}, rotate = 180] [fill={rgb, 255:red, 0; green, 0; blue, 0 }  ][line width=0.08]  [draw opacity=0] (8.93,-4.29) -- (0,0) -- (8.93,4.29) -- cycle    ;

% Text Node
\draw (146,357.4) node [anchor=north west][inner sep=0.75pt]    {$a$};
% Text Node
\draw (481,357.4) node [anchor=north west][inner sep=0.75pt]    {$b$};
% Text Node
\draw (311,466.4) node [anchor=north west][inner sep=0.75pt]    {$C_{0}$};
% Text Node
\draw (300,244.4) node [anchor=north west][inner sep=0.75pt]    {$C'_{0}$};
% Text Node
\draw (194,266.4) node [anchor=north west][inner sep=0.75pt]    {$U$};
% Text Node
\draw (170,375.4) node [anchor=north west][inner sep=0.75pt]    {$\nu '( p) =\tau ( p)$};
% Text Node
\draw (545.3,331.5) node [anchor=north west][inner sep=0.75pt]    {$\tau ( q)$};
% Text Node
\draw (466,407.4) node [anchor=north west][inner sep=0.75pt]    {$\nu '( q)$};

\end{tikzpicture}
}}

% 图 7.2：以用户提供的坐标稿为准；上方横轴的 0,1,2,3 等距。
\newcommand{\FigSevenTwo}{%
\begin{tikzpicture}[>=Stealth,x=.54pt,y=.54pt,yscale=-1,
    line width=.65pt,line cap=round,line join=round,font=\small]
  % 左上角的三个坐标方向。
  \draw (37,214)--(137,214) (47,124)--(47,224);
  \draw (130,209)--(137,214)--(130,219);
  \draw (42,131)--(47,124)--(52,131);
  \draw[->] (47,214)--(104.2,261.6);

  % 上方主坐标图。
  \draw (178.2,214)--(600.2,214);
  \draw (312,87.6)--(312,329.6);
  \draw (312,151)
    .. controls (324.2,169.6) and (351.2,184.6) .. (451.2,184.6)
    .. controls (551.2,184.6) and (551.2,240.6) .. (452.2,242.6);
  \draw (340,147)--(358.6,165.6);
  \draw (384.6,190.6)--(478.6,284.6);
  \draw (420.2,241.6) .. controls (382.2,241.6) and (330.2,259.6) .. (312.2,279.6);
  \draw[->] (467,152) .. controls (461.4,146.79) and (450.96,153.13) .. (449.2,172.6);

  % 由 1、2 的坐标间距确定 0、3：四点的横坐标依次相差 95.6。
  \foreach \x/\lab in {217.2/0,312.8/1,408.4/2,504.0/3}{
    \fill (\x,214) circle[radius=1.35pt];
    % 四个数字统一置于点的左下方；略靠近横轴，但字形仍完全位于轴下。
    \node[below left=2pt] at (\x,214) {$\lab$};}

  % 下方的球面同痕图。
  \draw (26.2,418.6)
    .. controls (27.2,403.6) and (38.4,385.2) .. (73.2,413.6)
    .. controls (108,442) and (136,450) .. (164.2,447.6)
    .. controls (192.4,445.2) and (205.2,456.6) .. (206.2,478.6)
    .. controls (207.2,500.6) and (162.2,507.6) .. (130.2,494.6)
    .. controls (98.2,481.6) and (96.2,473.6) .. (74.2,461.6)
    .. controls (52.2,449.6) and (25.2,433.6) .. (26.2,418.6)--cycle;
  \draw (463.2,457.6)
    .. controls (410.2,415.6) and (440.2,380.6) .. (459.07,377.85)
    .. controls (477.95,375.1) and (499.7,373.6) .. (523.2,379.6)
    .. controls (546.7,385.6) and (563.07,389.35) .. (583.14,414.48)
    .. controls (603.2,439.6) and (610.7,452.1) .. (611.2,458.6)
    .. controls (611.7,465.1) and (615.2,493.6) .. (593.2,498.6)
    .. controls (571.2,503.6) and (547.2,500.6) .. (534.2,496.6)
    .. controls (521.2,492.6) and (494.2,481.2) .. (474.2,469.2);
  \draw (191,452)
    .. controls (236.2,470.6) and (311.2,468.6) .. (372.2,467.6)
    .. controls (433.2,466.6) and (479.32,461.73) .. (507.2,462.6)
    .. controls (535.07,463.48) and (523.8,484) .. (514.8,486);
  \draw (242.2,489.6) .. controls (276.2,491.6) and (482.2,496.8) .. (506,489.2);
  \draw (201.2,489.6) .. controls (214.2,488.6) and (202.2,489.6) .. (215.2,489.6);
  \draw (175.2,470.6)
    .. controls (209.2,469.6) and (286.2,481.6) .. (362.2,480.6)
    .. controls (438.2,479.6) and (482.2,480.2) .. (538.2,466.6);

  % 变形方向箭头。
  \draw[<-] (214.59,481.39) .. controls (256.94,492.13) and (233.17,521.68) .. (271.2,524.6);
  \draw[<-] (188.2,507.53) .. controls (196.1,526.55) and (233.2,553.1) .. (271.2,524.6);
  \draw[<-] (152.92,437.86) .. controls (189.74,411.54) and (185.2,440.85) .. (224.2,411.6);
  \draw[<-] (225.98,452.44) .. controls (259.32,425.42) and (185.2,440.85) .. (224.2,411.6);

  \foreach \x/\y in {205.6/471.6,522.6/469.6,485.4/476}
    {\fill (\x,\y) circle[radius=1.35pt];}

  \node[anchor=west] at (474,138.4)
    {$H_1(1\times\mathbb R^{\lambda-1}\times0)$};
  \node[anchor=west] at (281,517.4) {$S_L=F_0(S_L)$};
  \node[anchor=west] at (216,387.4) {$F_1(S_L)$};
  \node[below=2pt] at (485.4,476) {$b$};
  \node[below left=2pt] at (205.6,471.6) {$a$};
  \node[anchor=west] at (120,186) {$(0,3)$};
  \node[above=2pt] at (47,124) {$\mathbb R^{\lambda-1}$};
  \node[below right=1pt] at (104.2,261.6) {$\mathbb R^{n-\lambda-1}$};
  \node[above=2pt] at (312,87.6) {$1\times\mathbb R^{\lambda-1}\times0$};
  \node[anchor=west] at (558,365) {$S_R$};
  \path[use as bounding box] (18,48) rectangle (670,565);
\end{tikzpicture}}

% 定理 8.1 指标 1 情形开头的临界点分层示意图。
\newcommand{\FigEightIndexOneSchematic}{%
\begin{tikzpicture}[x=1.35cm,y=1.0cm,line cap=round,line join=round]
  % 三个依次出现的水平流形；形状与图 3.1 保持一致。
  \draw[thick] (-3,-1.65)..controls(-3.16,-1.05)and(-2.86,-.52)..(-3,.02)
    ..controls(-3.14,.58)and(-2.88,1.12)..(-3,1.65);
  \draw[thick] (-1,-1.65)..controls(-.84,-1.05)and(-1.16,-.52)..(-1,.02)
    ..controls(-.86,.58)and(-1.14,1.12)..(-1,1.65);
  \draw[thick] (1,-1.65)..controls(.84,-1.05)and(1.16,-.52)..(1,.02)
    ..controls(.86,.58)and(1.14,1.12)..(1,1.65);

  % 各相邻水平之间的临界点列，指标依次为 1、2、3。
  \foreach \x in {-2,0,2}{
    \foreach \y in {-1.20,-.60,0,.60,1.20}{\fill (\x,\y) circle (1.65pt);}}

  \node[above=8pt] at (-3,1.65) {$V$};
  \node[above=8pt] at (-1,1.65) {$V_{1+}$};
  \node[above=8pt] at (1,1.65) {$V_{2+}$};
  \node[above=8pt] at (2,1.65) {等等};

  \node[below=10pt] at (-3,-1.65) {指标};
  \node[below=10pt] at (-2,-1.65) {$1$};
  \node[below=10pt] at (0,-1.65) {$2$};
  \node[below=10pt] at (2,-1.65) {$3$};
  \path[use as bounding box] (-3.85,-2.18) rectangle (2.70,2.18);
\end{tikzpicture}}
